\chapter{问答与对话}
\label{chap:nlp-question-answering-dialogue}

一份通知能够被准确复述，并不意味着系统已经能够回答围绕它提出的问题。用户问“我明天下午到可以吗”，系统还需知道咨询发生在哪一天、用户属于哪类参加者，以及哪些条件尚未满足。若下一轮用户更正了日期，系统必须撤掉旧要求；若查询失败，回复又不能把预期结果说成已经成立的事实。第\ref{chap:nlp-text-generation}章讨论怎样形成文本，本章增加两种约束：回答必须有可检查的依据，交流必须维护当前有效的目标。

本章的例子均为教学构造。共同通知发布于北京时间2025年4月15日，内容为：“机器学习读书会原定于4月17日举行，现改至4月18日（周五）14:00，地点为海棠楼201。校内师生可直接参加，校外来宾须在4月16日18:00前提交报名信息。”活动编号为\code{ML-0418}。通知没有提供实际报名人数、主讲人或报名成功记录；后文若需要这些数据，会另行声明教学表格或假设业务状态，不把它们归入这份通知。

围绕同一材料，可以提出不同任务。问答寻找并组织问题所需的事实；核验判断既有陈述与证据的关系；任务型对话把用户需求转成状态和业务动作；开放交流还需要区分知识、人物资料与情绪线索；协作对话则让要求在讨论中逐渐明确。它们能够共享编码、检索和生成部件，但不能共享一个含混的“回答质量”作为完成标准。

\section{证据问答}
\label{sec:nlp-dialogue-evidence}

证据问答的输入不只有问题$q$，还包括允许访问的资料、资料版本、查询时点和计算预算。输出可以是原文跨度、是非结论、列表、计算值或长回答，也可以是“给定材料不足以回答”。这里的不足是对可用证据的判断，不意味着世界上没有答案。连续文本、结构化数据和两者的混合需要不同的恢复过程，但都应保留从回答陈述回到来源的对应。

图\ref{fig:nlp-dialogue-evidence}把取得资料、形成答案与核查引用分开。问题同时要求时间和报名条件时，只找回时间通知属于覆盖失败；找到两个事实却把报名截止当成活动时间，属于组合错误；答案恰好正确却引用了不相关片段，仍是引用错误。多跳查询只在已有证据留下明确子问题时继续，不能用不断搜索掩盖已经耗尽的预算。

\begin{figure*}[htbp]
  \centering
  % Adapted from academic-drawing/templates/tikz_template.tex to the book styles.
% Main ports: south -> north. The dashed loop carries a new subquestion.
\begin{tikzpicture}[
  every node/.style={font=\figurefont\small,align=center},
  block/.style={book node,text width=64mm,minimum height=12mm},
  note/.style={text width=51mm,align=left,text=BookMuted},
  flow/.style={book arrow,rounded corners=1.2mm},
  retry/.style={flow,dashed,draw=BookAmber}
]
  \node[block,fill=FigureInput!12,draw=FigureInput] (question)
    at (-22mm,0) {问题与访问协议\\资料版本、时点、预算};
  \node[block,fill=FigureInput!12,draw=FigureInput] (candidates)
    at (-22mm,-23mm) {候选资料\\召回A、B，排除干扰C};
  \node[block,fill=FigureModel!12,draw=FigureModel] (evidence)
    at (-22mm,-46mm) {选定证据\\保留来源、对象与限定条件};
  \node[block,fill=FigureModel!12,draw=FigureModel] (answer)
    at (-22mm,-69mm) {答案陈述与来源标记\\时间引用A，截止条件引用B};
  \node[block,fill=FigureOutput!12,draw=FigureOutput] (check)
    at (-22mm,-92mm) {引用核验与交付\\支持则保留，不足则限定或弃答};

  \draw[flow] (question.south) -- (candidates.north);
  \draw[flow] (candidates.south) -- (evidence.north);
  \draw[flow] (evidence.south) -- (answer.north);
  \draw[flow] (answer.south) -- (check.north);
  \node[note] at (48mm,-23mm) {检索遗漏\\只召回B，缺活动时间};
  \node[note] at (48mm,-46mm) {组合错误\\把报名截止当成活动时间};
  \node[note] at (48mm,-69mm) {无支持引用\\引用A解释校外报名条件};
  \node[note] at (48mm,-92mm) {独立验收\\答案、证据覆盖、引用支持};
  \draw[retry] (evidence.west) -- (-66mm,-46mm)
    -- (-66mm,0) -- (question.west);
  \node[font=\figurefont\footnotesize,text=BookAmber,rotate=90]
    at (-72mm,-23mm) {未决子问题：预算内再查询};
\end{tikzpicture}

  \caption{证据取得、答案形成与引用核验。实线表示当前问题的数据流，虚线表示预算内由未决子问题触发的再查询。教学资料固定为时间地点片段A、报名条件片段B和干扰片段C；答案是否正确、所需证据是否齐全、引用是否支持相应陈述分别验收。}
  \label{fig:nlp-dialogue-evidence}
\end{figure*}

\subsection{连续文本问答}
\label{sec:nlp-dialogue-text-qa}

给定完整通知，“何时举行”可以恢复为“2025年4月18日14:00”；“谁可直接参加”要求类别列表“校内师生”，而不是实际出席名单；“主讲人是谁”在材料中没有答案。SQuAD 2.0把这种证据相对的不可回答性纳入阅读任务，要求模型区别“出现一个类型相似的短语”和“短语确实回答问题”\scite{rajpurkar2018squad2}

从资料集合回答还多出候选取得的环节。系统先切分资料并保存文档编号、原文偏移和版本，再按问题召回片段，必要时重排，组织阅读上下文，恢复答案并检查支持。多跳问答还可能用一条事实中的对象触发下一次查询；HotpotQA提供支持事实标注，使答案和证据组合可以分别评价（见本章末文献说明）。长文实验应从原始长材料开始计算覆盖率，不能只给模型已经包含答案的窗口，再把阅读准确率称作端到端准确率。

\subsubsection{BERT跨度阅读与弃答}
\label{sec:nlp-dialogue-bert-span}

若答案是材料中的连续文字，最直接的预测对象是起止位置。把问题和材料组成\code{[CLS] 问题 [SEP] 材料 [SEP]}，由BERT得到位置表示$\bm{h}_i\in\mathbb{R}^{d}$。编码机制见第\ref{chap:transformer}章；问答任务另外学习起点向量$\bm{w}_{\mathrm{s}}$和终点向量$\bm{w}_{\mathrm{e}}$。材料位置与空答案位置参与归一化，问题位置不作为可返回跨度：
\begin{equation}
  p_{\mathrm{s}}(i)=
  \frac{\exp(\bm{w}_{\mathrm{s}}^\top\bm{h}_i)}
       {\sum_{j\in I}\exp(\bm{w}_{\mathrm{s}}^\top\bm{h}_j)},
  \qquad
  L=-\log p_{\mathrm{s}}(i^*)-\log p_{\mathrm{e}}(j^*).
  \label{eq:nlp-dialogue-span}
\end{equation}
其中$I$是允许位置集合，终点分布$p_{\mathrm{e}}$以同样方式计算。训练样本包含问题、材料、正确字符区间或无答案标签；分词映射将区间转为$(i^*,j^*)$，无答案样本将两端都指向\code{[CLS]}。任务损失同时更新编码器和两个位置预测头\scite{devlin2019}

解码时不能分别选择两个最大概率位置：这样可能得到终点在起点之前的结果。应在$i\leq j$且$j-i+1\leq\ell_{\max}$的合法跨度中最大化起止分数之和。设最佳跨度的未归一化分数为$s_{\mathrm{span}}$，空位置两端之和为$s_{\mathrm{null}}$，当$s_{\mathrm{span}}-s_{\mathrm{null}}>\tau$时返回跨度，否则弃答。阈值$\tau$和最大长度在独立开发集上选定；阈值不是用户语气中的“我很确定”。

用压缩后的材料词元“读书会／改至／4月18日／14:00／地点／海棠楼201／校内师生／可直接参加”演算。对时间问题，假设位置3和4的起点概率为$0.60,0.10$，终点概率为$0.15,0.70$，则$(3,4)$的乘积分数为$0.42$，两个单词元候选分别为$0.09,0.07$。若空答案的乘积分数为$0.02$，相同归一化下分数差为$\log(0.42/0.02)=3.045$，在教学阈值$\tau=1$时返回“4月18日14:00”。位置到字符的映射再补回原文空白或标点，而不是把模型子词拼接当作原文。

同一材料对“谁可直接参加”可把位置7作为单跨度，假设其乘积分数$0.56$，空答案分数$0.04$，差值$\log14=2.639$，仍接受。要演示“未提供地点”，必须另给删去地点字段的教学材料，而不能声称共同通知没有地点。该材料上若最佳错误跨度分数$0.03$、空分数$0.30$，差值$-2.303$，就应弃答。完整通知上的“主讲人是谁”也属于应弃答的问题。

若问题要求“校内学生和教师”这类不连续列表，需要多跨度标签或列表生成；“校外来宾能否直接参加”需要是非或条件答案预测。单跨度头不会自动获得这些输出类型。长材料可以滑动分窗，对每窗保留偏移、最佳非空候选与空分数，按预先规定的跨窗评分汇总并去重。每个窗口内部单独归一化的概率不可无条件直接比较；可用统一未归一化分数加开发集校准。跨窗分散的答案不能靠单窗跨度解码恢复，还需增加组合机制。

\subsubsection{BiDAF问文交互阅读}
\label{sec:nlp-dialogue-bidaf}

同样的跨度监督也能接到另一种问文交互表示上。双向注意流（bidirectional attention flow, BiDAF）让每个材料词保留自己的问题相关表示，而不是先把整段压成一个向量。原模型将预训练词向量与字符卷积表示连接，经highway网络和双向LSTM得到材料向量$\bm{h}_i$与问题向量$\bm{u}_j$；随后学习相似矩阵
\begin{equation}
  \bm{S}_{i,j}=\bm{w}^{\top}
  [\bm{h}_i;\bm{u}_j;\bm{h}_i\odot\bm{u}_j].
  \label{eq:nlp-dialogue-bidaf}
\end{equation}
其中$\odot$表示逐元素乘积。每一行做softmax得到材料到问题的注意力$\alpha_{i,j}$，从而$\widetilde{\bm{u}}_i=\sum_j\alpha_{i,j}\bm{u}_j$。反方向先对每行取最大值，再在材料位置间做softmax，得到$\beta_i$及材料摘要$\widetilde{\bm{h}}=\sum_i\beta_i\bm{h}_i$。送入后续建模层的是$[\bm{h}_i;\widetilde{\bm{u}}_i;\bm{h}_i\odot\widetilde{\bm{u}}_i;\bm{h}_i\odot\widetilde{\bm{h}}]$，并非某一个注意力最大的词\scite{nlp-dialogue-seo2017}

以两行、两列的教学子矩阵
$\bm{S}=\left[\begin{smallmatrix}2&0\\0&1\end{smallmatrix}\right]$
说明汇总。两行权重分别为$(0.881,0.119)$和$(0.269,0.731)$；行最大值$(2,1)$产生反向权重$(0.731,0.269)$。这些权重决定送往跨度头的表示，不直接等于答案概率。后续双向LSTM建模材料关系，起点头使用交互表示和建模表示，终点头还使用进一步的循环编码；两端位置交叉熵训练所有任务参数，合法跨度最大乘积恢复答案。

固定上一小节的通知和时间答案，BERT与BiDAF改变的是任务表示的形成方式，位置标签和字符恢复契约可以保持一致。另设一组歧义材料：“甲校的林宁主持读书会，乙校的林宁主持答疑会。”问题为“谁主持乙校答疑会”。仅匹配“林宁”会留下两个同名位置；学校和活动关系共同决定应返回第二个提及的字符区间。抽出的字符串相同仍不足以证明定位正确。原始BiDAF的阅读设置不包含本书的全部弃答、长文分窗和列表协议；增设空答案头时需要对应负例，不能把这些能力归因于名称中的“双向”。

\subsubsection{DPR检索与阅读}
\label{sec:nlp-dialogue-dpr}

跨度阅读假定材料已经在手。现在把共同通知拆成互补片段A和B：A保留活动编号、改期、时间和地点，B保留同一编号及校内、校外参加条件；另有无关的教室借用说明C。问题是“读书会什么时候举行，校外来宾报名何时截止”。检索器需要让A、B都进入候选，读者才可能恢复完整回答。片段保存$(\text{文档编号},\text{版本},\text{起止偏移})$，切分时保留编号和标题，避免把脱离对象的“18:00前”当成独立事实。

无需问题标注的词汇基线是BM25。它按词在资料集合中的稀有程度加权，对片段内重复词计数作饱和处理，再校正片段长度。采用下式的非负逆文档频率约定：
\begin{equation}
\begin{aligned}
 s_{\mathrm{BM25}}(q,z)
 &=\sum_{w\in q}
 \log\left(1+\frac{n-n_w+0.5}{n_w+0.5}\right)
 \frac{f_{w,z}(k_1+1)}
      {f_{w,z}+k_1(1-b+b|z|/\bar{\ell})}.
\end{aligned}
\label{eq:nlp-dialogue-bm25}
\end{equation}
其中$n$为片段数，$n_w$为含词$w$的片段数，$f_{w,z}$为词频，$|z|$为片段词元数，$\bar{\ell}$为平均长度。索引统计来自固定资料快照；$k_1,b$是用开发查询选择的超参数，不是逐条问题学习出的参数。分词、停用词及同义词处理也是实验设定的一部分。

为把数值算清，表\ref{tab:nlp-dialogue-retrieval}直接规定一个教学索引的统计，查询只保留“读书会、报名”，取$k_1=1.2,b=0.75$。表中的长度是教学分词统计，不等于上文中文片段的字符数。$n=3,\bar{\ell}=40/3$，两个词的逆文档频率为$\log1.6=0.470$和$\log(8/3)=0.981$。例如B的分数为$0.470\times2.2/2.65+0.981\times4.4/3.65=1.573$。取前两项便得到B、A；只取第一项仍缺活动时间。

\begin{table}[htbp]
  \centering
  \begin{tabular}{lrrrr}
    \toprule
    片段 & 长度 & 读书会词频 & 报名词频 & BM25分数\\
    \midrule
    A：时间地点 & 10 & 1 & 0 & 0.524\\
    B：参加条件 & 20 & 1 & 2 & 1.573\\
    C：借用说明 & 10 & 0 & 0 & 0\\
    \bottomrule
  \end{tabular}
  \caption{固定教学索引上的词汇召回。片段A、B互补，排序靠前和证据覆盖完整不是同一个条件。}
  \label{tab:nlp-dialogue-retrieval}
\end{table}

若用户问“校外怎么参加”，字面上可能没有“报名”。稠密片段检索（dense passage retrieval, DPR）学习问题与片段的向量，使相关但用词不同的文本也能接近。分别用问题编码器$E_q$和片段编码器$E_z$取得$\bm{q},\bm{z}\in\mathbb{R}^d$，分数为$s(q,z)=\bm{q}^{\top}\bm{z}$。训练样本包含问题、正确片段$z^+$与负片段集合$Z^-$，最小化
\begin{equation}
 L_{\mathrm{DPR}}=
 -\log\frac{\exp s(q,z^+)}
 {\exp s(q,z^+)+\sum_{z^-\in Z^-}\exp s(q,z^-)}.
 \label{eq:nlp-dialogue-dpr}
\end{equation}
原论文用两个预训练BERT初始化编码器，比较随机负例、BM25难负例及批内其他问题的正确片段；批内复用可在一次编码后形成多个对比项。正例来自证据标注，或以答案匹配构造的弱监督。后者可能把同名但无关的片段误当正例；其他问题的正片段也不一定对当前问题为负，构造多跳数据时尤其要检查假负例（文献说明列出DPR原文）。

假设训练后对本题有$\bm{q}=(1,1)$，三个片段向量为$\bm{z}_A=(2,0)$、$\bm{z}_B=(0,3)$、$\bm{z}_C=(-1,0)$，内积依次为$2,3,-1$。对“报名截止”子问题以B为正例时，归一化正例概率为$e^3/(e^2+e^3+e^{-1})=0.721$，损失为$0.327$。这里的二维向量仅用来检查评分，不冒充真实训练所得表示。对需要两项事实的完整问题，A也应作为有效证据，不能因为本轮选择B为一个训练正例就把A永久标成负例。

使用时离线编码全部片段并建立内积索引，在线只编码问题，取$k$个候选。若再接交叉编码器，它将每个问题与候选联合编码，以相关性标签训练评分头，只能改变已召回候选的次序，不能补回集合之外的证据。索引近似搜索的遗漏也应与表示错误分开。编码器更新后，旧片段向量与新问题向量可能不再处于同一空间，因此需要记录索引和编码器版本。

取$k=2$得到B、A后，将每个片段送入跨度阅读器：A返回“4月18日（周五）14:00”，B返回“4月16日18:00前”。再用问题对应的字段模板恢复“活动于2025年4月18日14:00在海棠楼201举行；校外来宾须在4月16日18:00前提交报名信息”。时间地点引用A，报名条件引用B。第二分句中的“前”不能丢掉；它规定严格早于18:00。只用A引用整段，即使文字正确也缺报名依据。与直接长上下文比较时，固定资料快照、可见总词元、候选数和生成预算，并分别报告召回覆盖、阅读正确与引用支持，不能把三者压成一个检索分数。

\subsubsection{RAG检索条件生成}
\label{sec:nlp-dialogue-rag}

上一条路线分别读出字段，再用模板组合。若答案需要较自由的表达，可以让生成器在检索文档条件下预测答案。检索增强生成（retrieval-augmented generation, RAG）原模型把文档看作隐变量：不预先指定唯一正确文档，而把各文档条件下的答案概率加权求和。本小节特指这套概率模型，不把所有“检索后接语言模型”的系统都视作同一个算法。

设检索到的$k$个片段为$Z_k$，$\pi_z=p_\eta(z\mid q)$为候选内归一化的检索概率，生成器给出$p_\theta(y_i\mid q,z,y_{<i})$。序列级与词元级边际化分别为
\begin{equation}
\begin{aligned}
 p_{\mathrm{seq}}(y\mid q)
 &=\sum_{z\in Z_k}\pi_z
      \prod_{i=1}^{m}p_\theta(y_i\mid q,z,y_{<i}),\\
 p_{\mathrm{tok}}(y\mid q)
 &=\prod_{i=1}^{m}
      \sum_{z\in Z_k}\pi_z p_\theta(y_i\mid q,z,y_{<i}).
\end{aligned}
\label{eq:nlp-dialogue-rag}
\end{equation}
序列级模型在一次完整答案条件概率内使用同一文档，再对文档求和；词元级模型在每个位置都重新求和，因此允许不同位置依赖不同文档。后者不意味着每生成一个词就重新发出检索查询，原模型的候选集合和检索先验仍来自输入问题。

表\ref{tab:nlp-dialogue-rag}给出完整教学计算。为了突出求和与求积的位置，只比较两个双词元候选，省略两者相同的结束符概率；这些数值是给定各自前缀时的条件概率，不是文献实验。取$\pi_A=0.6,\pi_B=0.4$，候选甲为“周五／14:00”，候选乙为“周四／14:00”。

\begin{table}[htbp]
  \centering
  \begin{tabular}{llrr}
    \toprule
    候选 & 条件片段 & 第1词元概率 & 第2词元概率\\
    \midrule
    甲 & A & 0.8 & 0.2\\
    甲 & B & 0.2 & 0.9\\
    乙 & A & 0.1 & 0.5\\
    乙 & B & 0.4 & 0.4\\
    \bottomrule
  \end{tabular}
  \caption{RAG的教学条件概率。第2列表示生成时访问哪个片段；第4列还以该候选的第1词元为条件。}
  \label{tab:nlp-dialogue-rag}
\end{table}

序列级甲得分为$0.6(0.8)(0.2)+0.4(0.2)(0.9)=0.168$，乙为$0.6(0.1)(0.5)+0.4(0.4)(0.4)=0.094$。词元级甲得分为$(0.6\times0.8+0.4\times0.2)(0.6\times0.2+0.4\times0.9)=0.2688$，乙为$(0.6\times0.1+0.4\times0.4)(0.6\times0.5+0.4\times0.4)=0.1012$。两种模型都在这两个候选中选甲，但它们不是用同一个分数比较。甲在词元级模型下的第一位置文档后验为$(0.857,0.143)$，第二位置为$(0.25,0.75)$；后验改变来自生成概率，而非新的检索调用。

训练用问题与参考答案对，最小化相应边际概率的负对数。原RAG以DPR初始化检索器、以BART初始化生成器，在任务微调时固定文档编码器和索引，更新问题编码器与生成器；没有在该微调目标中要求逐题文档标签，不等于初始化也没有使用检索监督。词元级分布可直接接束搜索；序列级则先在各文档条件下产生候选，再对候选在各文档下重评分求和，或明确采用省略未生成候选项的近似。这些机制见文献说明中的RAG原文。

概率最高仍不等于证据支持。例如B只陈述报名条件，却可能因参数知识给“14:00”较高概率。要输出“周五14:00〔A〕”，需另外规定来源标记的输出语法，要求每个标记解析为已检索片段，并检查A实际支持该时间；模型的文档后验本身不是有效引用。对于多跳问题，可以在本书协议中把未解决子问题加入查询队列，记录已见来源，达到证据充分、无新资料或预算上限时停止。这是检索控制扩展，应分别报告检索缺失、错配和组合错误，不能归入式\eqref{eq:nlp-dialogue-rag}自动保证的能力。

\subsubsection{FiD多片段融合生成}
\label{sec:nlp-dialogue-fid}

RAG的每个文档条件生成分支只直接访问一个片段。若回答必须同时用到A的活动时间和B的报名截止，可改用解码器融合（fusion-in-decoder, FiD）：每个问题与片段对独立编码，把全部编码结果交给同一个解码器。构造输入“question: 问题 title: 标题 context: 片段”，得到
\begin{equation}
 \bm{H}_j=E_\theta(q,z_j),\qquad
 p(y\mid q,Z_k)=
 \prod_i p_\theta(y_i\mid y_{<i},[\bm{H}_1;\ldots;\bm{H}_k]).
 \label{eq:nlp-dialogue-fid}
\end{equation}
这里连接的是记忆序列，解码器在每一步对所有片段表示做交叉注意力。并没有给各文档的完整答案概率再乘$\pi_z$求和。原方法以预训练的序列到序列模型初始化，检索结果作为输入，用参考答案的词元负对数似然训练阅读生成部分；片段检索可用BM25或DPR，原文主实验使用T5并采用贪心解码（见文献说明）。

固定A、B两个片段，问题为“活动时间与校外报名截止分别是什么”。只给A的阅读器最多提供活动时间，只给B的阅读器最多提供报名截止；FiD解码器同时可见两者，能够生成“4月18日14:00举行，校外来宾应在4月16日18:00前提交报名信息”。输出仍要解析成两个陈述，再将时间对应A、截止条件对应B。若生成器把两个日期交换，两个片段都被访问过也不能为答案辩护。

设每个问题—片段输入长度为$\ell$，片段数为$k$，暂忽略表示宽度和层数，独立编码的自注意力配对量约为$k\ell^2$，直接连接长上下文约为$(k\ell)^2$。在$k=2,\ell=100$的教学设置下为20000与40000对位置；两者的解码交叉注意力仍需访问约200个输入位置，且FiD重复编码了问题。这只是注意力项的比较，不是总运行时间保证。长上下文编码允许片段在编码阶段交互，FiD将这种结合推迟到解码阶段。固定片段、总预算和资料版本后，才有意义比较这种位置安排怎样影响覆盖、成本和组合错误。

\subsection{结构化证据问答}
\label{sec:nlp-dialogue-structured-qa}

问题的依据若已经是表格、数据库或图记录，回答可以先执行明确的选择与计算，再把结果表达成文字。其输入除了问题，还包括模式、允许访问的数据快照和只读权限；其输出需要保存查询或选择结果，而不只有一句自然语言答案。第\ref{chap:nlp-text-analysis}章介绍从语言到形式表达的映射，这里关注执行结果怎样成为答案证据。

\subsubsection{SQL与图查询的执行式问答}
\label{sec:nlp-dialogue-execution-qa}

设有额外构造的报名表\code{registration}，字段为活动编号、报名记录编号、参加者编号、身份类别、提交时刻和随行人数。四条记录为：R1，校外，4月16日17:30，2人；R2，校外，4月16日18:00，1人；R3，校内，4月17日09:00，1人；R4，校外，4月16日16:00，3人。它们都属于\code{ML-0418}，时区统一为北京时间，记录编号唯一。本例只问“按时提交的校外报名有几条，合计几人”，不把提交记录等同于已批准资格。

模式链接将“校外”对应身份字段，“按时”对应通知中严格早于截止时刻的条件，“几条”与“几人”分别对应记录数和人数之和。语义解析器可由问题—查询对监督训练，也可在给定模式和示例下生成候选查询；解析后检查字段、类型、只读性和条件，再绑定参数执行：
\begin{codeblock}[language=SQL,label={code:nlp-dialogue-sql},unbreakable]{报名统计的只读查询}
SELECT COUNT(*) AS records,
       SUM(party_size) AS people
FROM registration
WHERE event_id = :event
  AND visitor_type = :external
  AND submitted_at < :deadline;
\end{codeblock}
绑定的活动为\code{ML-0418}，截止为2025年4月16日18:00。筛选留下R1、R4，结果为2条、5人。回复应为“按给定快照，截止前提交的校外报名有2条，登记人数合计5人”，并保留行编号R1、R4。若把条件误写成$\leq$，会多计R2，得到3条、6人；查询执行无错误，却没有回答原问题。

图数据可把活动、报名记录和参加者表示为节点，以“属于活动”“提交者”等边连接。问题到图路径的映射仍需绑定活动身份和时间条件，再计数通过筛选的报名节点；跨边连接时保存节点及边编号。若一个人关联两条报名记录，计数报名节点与计数不同参加者就不同。SQL的多表连接也可能重复行，因此要在问题要求人数去重时明确采用参加者标识，而不能任意增加\code{DISTINCT}让数字看起来合理。

恢复过程应分别检查解析、执行、组合和表达。两条语法不同但等价的查询可以都正确；两条不等价查询也可能在当前表上偶然同值。用截止恰为18:00、同名不同编号、重复连接等反例数据复核语义，比只比较一次执行结果更能暴露错误。空结果也不总是零：\code{COUNT}可得0，某些数据库中的空集\code{SUM}为\code{NULL}，需要由答案协议决定表达为“无符合记录”还是数值零。

\subsubsection{TAPAS单元格选择与聚合}
\label{sec:nlp-dialogue-tapas}

当模式较简单时，可以不生成SQL，而直接指出哪些单元格参与回答。TAPAS把问题与按行展开的表格输入BERT式编码器，除词元和位置外加入行、列、数值排序等表示；对一个单元格内的词元logit取平均，再经sigmoid得到选择概率$p_c$。聚合头从全局表示预测\code{NONE}、\code{COUNT}、\code{SUM}或\code{AVERAGE}，推断时选择概率超过阈值的单元格，执行所选算子。原模型还采用单列选择的归纳偏置，不能据此直接处理任意跨表程序（见文献说明）。

监督强度需要说清。若答案文字能够映射回单元格，训练可使用列选择损失、单元格二分类损失和无聚合标签。若只知道标量答案而没有操作过程，TAPAS用可微近似计算软计数$\sum_c p_c$、软求和$\sum_c p_cv_c$等，再结合聚合概率得到软答案，以数值误差和聚合约束训练；原文使用Huber型误差。它不是在每条训练样本上都拥有人工标注的SQL，也不能把软求和称作最终离散执行结果。

固定上一小节的表，问合计人数时假设四个人数单元格选择概率为$(0.95,0.10,0.05,0.90)$。阈值$0.5$留下R1的2和R4的3，\code{SUM}返回5人；训练时软求和却为$0.95\times2+0.10\times1+0.05\times1+0.90\times3=4.75$，与目标5尚有误差。问有几条时可以选择记录编号列中的R1、R4并执行\code{COUNT}得到2；问哪些记录则用\code{NONE}恢复编号列表。列选择和聚合算子都必须由问题决定，不能对同一组数值不加区分地操作。

若两条记录的参加者都叫“林宁”，只输出名字无法区分身份，需在协议中保留记录编号。若表把一行的量记为“2组”、另一行为“3人”，即使数值相加等于5，也不能表达为5人。弱答案监督还可能让错误选择集偶然得到同一合计，因此评价除答案外应检查单元格、筛选条件和单位。需要跨列比较、任意连接或多步组合时，执行式查询通常给出更明确的操作空间；这不意味着生成查询就免除了语义错误。

\subsection{混合证据问答}
\label{sec:nlp-dialogue-hybrid-qa}

有些条件只在正文中说明，数值却在表格里。此时仅找出数字不足以形成答案，还需把对象、时间、单位和比较范围对齐。本节从已可用的文本和结构化表格出发，不处理页面图像、表格识别或版面恢复。表\ref{tab:nlp-dialogue-operands}是独立的教学统计，与共同通知的报名记录无关。附带说明为“比较甲校读书会的年度到场人数；两年口径相同，均按去重人数统计，不含旁听记录”。

\begin{table}[htbp]
  \centering
  \begin{tabular}{llll}
    \toprule
    单元格编号 & 对象 & 年份 & 到场人数\\
    \midrule
    T1-R1-C3 & 甲校读书会 & 2024 & 120人\\
    T1-R2-C3 & 甲校读书会 & 2025 & 150人\\
    T1-R3-C3 & 乙校读书会 & 2025 & 180人\\
    \bottomrule
  \end{tabular}
  \caption{增长率问题的教学操作数。文本说明限定甲校、年度和去重口径，表格提供数值及年份；乙校一行只用作对象对齐的反例。}
  \label{tab:nlp-dialogue-operands}
\end{table}

\subsubsection{TAGOP标签选择与符号计算}
\label{sec:nlp-dialogue-tagop}

TAT-QA提出的TAGOP把混合问答分成证据选择、算子选择、必要的操作数排序和尺度恢复。问题、按行展开的表格与相关正文经编码后，每个子词预测\code{I}或\code{O}标签，表示是否进入证据候选；单元格或原词只要有子词标为\code{I}就作为正候选，正文中连续的正词合并成跨度。这个IO接口不同于BIO实体边界标注，恢复时必须保留表格坐标与原文偏移（见文献说明中的TAT-QA原文）。

算子分类头从问题条件的全局表示预测预定义操作，包括单跨度、单元格、多项返回、求和、计数、平均、乘法、除法、差额和变化率。对差额、除法和变化率，前两个高分数值还需经过顺序判定，因为$150-120$与$120-150$不是同一答案。尺度头结合全局、表格和正文汇总表示预测无尺度、千、百万、十亿或百分比。监督来自答案与推导标注恢复的选择标签、算子、操作数顺序和答案尺度；训练将这些分类损失相加。原文以首次匹配定位重复证据的做法可能产生监督歧义，教材例子保留唯一单元格编号以便检查。

问题“甲校2025年到场人数比2024年增长百分之几”应选T1-R2-C3的150人和T1-R1-C3的120人，并以附带文本确认比较范围。假设两个数值的选择置信分别为$0.93,0.89$，乙校180人的置信为$0.08$，算子选择\code{Change ratio}，顺序确定为当年数、上年数。执行得到
\begin{equation}
 r=\frac{150\ \text{人}-120\ \text{人}}{120\ \text{人}}
   =0.25,\qquad \text{百分比答案}=25\%.
 \label{eq:nlp-dialogue-growth}
\end{equation}
人数单位在比值中消去；输出的25与百分号必须一起恢复，不能把0.25直接写成0.25\%。操作记录应同时保存两个单元格、文本限定和尺度。若问题改问增加多少人，算子应改为差额，结果为30人，而不是复用25\%。

乙校180人即使与甲校120人构成可执行变化率，也回答了错误对象。若上一年写成0.12千人，应先显式规范为120人；若基数为0，增长率无定义，应报告无法按该公式计算。TAGOP的一次预定义操作适合其操作集合覆盖的题目；需要“两个部门先合计，再减去旁听人数，然后求比值”的题目不能仅靠增加一个输出尺度完成。

\subsubsection{FinQANet程序生成与执行}
\label{sec:nlp-dialogue-finqa}

复杂计算可以把中间结果明确写出来。FinQANet将表格行转成带行列含义的句子，与正文一起组成事实候选，先用问题—事实分类器检索相关项，按报告原顺序组织选中事实，再生成限定语言中的计算程序。检索器以支持事实标签训练；生成器以参考程序监督，而不是只以最终数字判断每一步是否正确（见文献说明中的FinQA原文）。

程序词元有三个来源：输入中的数字或表格行名，限定语言中的函数、括号和常数，以及表示既有步骤结果的记忆标记，例如\code{\#0}。原模型用预训练编码器表示事实和问题，以LSTM解码；解码状态分别注意输入与已生成程序，再为可选词元评分。每完成一个操作的右括号，就用当前上下文更新对应步骤的记忆表示，后续步骤可选择该标记。这里的学习式记忆表示帮助生成引用，实际算术值仍由执行器计算，不应将两者混为一谈。语法掩码可以阻止参数位输出非法函数，却不会判断当前数字是否来自正确年份。

固定表\ref{tab:nlp-dialogue-operands}和相同文本，不改变证据预算，教学程序为：
\begin{codeblock}[language=Python,label={code:nlp-dialogue-finqa}]{从差额到增长率的限定程序}
subtract(150, 120)   # step 0: 30 persons
divide(#0, 120)      # step 1: 0.25
\end{codeblock}
其中150绑定T1-R2-C3，两个120都绑定T1-R1-C3，\code{\#0}只允许引用已完成的第0步。此块展示的是限定语言，不是可直接执行的Python语法。执行器先产生30人，再以120人为分母得到无量纲比值0.25；答案格式化为25\%，并附上文本限定和单元格来源。若选择\code{divide(120,\#0)}，程序语法正确但结果4不符合问题；若把两个120替换成碰巧同值的其他部门记录，执行答案也可能不变，但证据绑定已经错误。

相对于TAGOP内置的变化率算子，FinQANet把同一计算展开为两个可组合步骤，更容易表示其限定语言支持的多步程序，也引入步骤选择、引用和终止错误。直接答案生成可以在固定证据下作为基线，但没有程序可执行检查。程序准确率应允许预先规定的等价程序，例如在适用数域内等价的重排，同时另报执行答案准确率；两者都不能省略对象、单位与时间检查。通用程序生成、测试和修订循环由第\ref{chap:nlp-solving-and-execution}章展开。

\section{主张与回答核验}
\label{sec:nlp-dialogue-verification}

问答生成的是新答案，核验检查的是已经提出的主张。核验输入应包含待查文本、允许使用的知识源、时点和证据访问协议，输出则需要指出哪个陈述受到支持、被反驳或尚无足够信息。一次整体评分无法告诉读者该保留哪句话；可操作的核验需要保留主张范围、依据和修改理由。

\subsection{主张拆解与证据定位}
\label{sec:nlp-dialogue-claims}

“读书会4月18日在海棠楼201举行，现场有150人”同时包含日期、地点和人数。不能因为前两项正确就整体判真，也不能因为人数缺依据就丢掉已经获得的时间地点。拆解的任务是形成可以分别检索和判断的命题，且不改变否定、数量词、对象与时间条件；不是机械地按句号切分。

\subsubsection{FActScore原子事实分解}
\label{sec:nlp-dialogue-factscore}

FActScore把长回答分成原子事实，即每条只承载一个可独立检查的信息点，再计算指定知识源支持的比例。其自动估计器用语言模型形成事实单元，检索有关资料，然后逐条判定支持；原研究在人物传记中评价事实精确性，不把这个比例定义为完整性、信息量或全部回答质量。自动分解以提示和示例驱动，论文的人类标注流程还允许修订机器提出的单元，不能把自动列表当成无需检查的真值（见文献说明）。

处理上面的回答时，先将“读书会”绑定\code{ML-0418}，把隐含年份恢复为2025年，得到三条记录：活动在4月18日举行；活动地点为海棠楼201；活动现场人数为150人。每条记录保留原句跨度及补入背景的来源，查询分别包含活动编号与待查属性。为演示“只有人数错误”，另设一份明确的教学签到统计S：“ML-0418实际到场人数为120人，统计时点为活动结束后。”通知支持前两条，S反驳第三条，于是未加额外惩罚的单回答支持比例为
\begin{equation}
 \operatorname{FP}(y;C)=
 \frac{\sum_{a\in A_y}\mathbb{I}[C\text{支持}a]}{|A_y|}
 =\frac{2}{3}.
 \label{eq:nlp-dialogue-factscore}
\end{equation}
其中$A_y$是本次完整拆解的原子事实集合，$C$是规定的知识源。若没有签到统计S，人数条目只能标为未获支持，不能仅凭通知未提人数就声称150为假；两种情形的支持比例相同，核验原因却不同。

代词恢复也会影响分母。例如“林宁主持了活动。他来自甲校”必须把第二句的“他”绑定同一人物，否则可能对另一名林宁查出支持。分解器若遗漏人数条目，分数会虚升至$2/2$；若把“至少150人”改成“150人”，又改变了原意。计算比例前应检查覆盖和含义保持，空事实集合另报不可评分，不能把它默认为满分。按模型汇总时还应报告回答率和每条回答的事实数量，防止用少说或不说掩盖信息缺失。

\subsubsection{规则分解与结构抽取}
\label{sec:nlp-dialogue-rule-claims}

在通知和业务回复中，事实句式常较稳定，可以用规则建立可检查的拆解基线。规则识别活动主语、时间、地点、人数与限定词，或者复用第\ref{chap:nlp-information-extraction}章的实体、关系和事件记录，再把各字段恢复为命题。规则无需从样本估计全部参数，但需要人工规定适用句式，并以独立标注样本检查覆盖；使用学习式抽取器时，其参数则来自实体、关系或事件监督。

对同一回答，可以先恢复事件记录$(\text{ML-0418},\text{日期}=4/18,\text{地点}=\text{海棠楼201},\text{人数}=150)$，再输出与上一小节相同的三个主张，并保留每个字段的字符区间。遇到“原定4月17日，现改4月18日”，规则必须生成“原计划日期”和“当前日期”两个不同关系；简单抽取两个日期再任选其一，会把时间更新误当来源冲突。遇到“不是所有校外来宾都可直接参加”，还要保留量词和否定作用域，不能转换为“所有校外来宾都不可参加”。

规则与模型的比较应固定同一回答和参考命题集合，分别计算遗漏、额外主张和含义改变。某个结构记录通过类型检查，只说明它能被解析；记录是否表达了完整命题仍要检查。例如“校外来宾，报名，4月16日”漏了18:00和严格的“前”，不能作为与原通知等价的核验单元。资料不足或来源冲突应留到证据判断阶段，不能在抽取时删掉难以验证的字段。

\subsection{支持、反驳与信息不足}
\label{sec:nlp-dialogue-evidence-relation}

给定一个主张和所取得的证据，FEVER式协议区分支持、反驳和信息不足：证据能推出主张、证据能推出与其不相容的内容，或当前证据不足以作出前两种判断。该接口与第\ref{chap:nlp-classification-matching-organization}章的句间推断相连，但还增加了证据是否取得、来源是否适用的问题。无关段落不是反证，同一句话被多次转载也不会自动增强其事实依据。

\subsubsection{可分解注意力的证据判定}
\label{sec:nlp-dialogue-decomposable}

FEVER原始流水线先用词汇检索取得文档，再选择证据句，将选择的句子连接后交给可分解注意力模型，而不是对每句三分类后简单多数投票。关系模型以有证据的主张和关系标签监督训练；信息不足样本还需要构造输入证据，原文比较从相关检索结果取句和随机取句两种方式。这种训练分布差异会影响实际检索失败时的行为\scite{nlp-intro-thorne2018}

可分解注意力把计算拆成对齐、比较和聚合。设证据词表示为$\bm{e}_i$，主张词表示为$\bm{c}_j$，学习前馈变换$F$，令$s_{i,j}=F(\bm{e}_i)^\top F(\bm{c}_j)$。按行归一化得到每个证据词对应的主张摘要$\widetilde{\bm{c}}_i$，按列归一化得到每个主张词对应的证据摘要$\widetilde{\bm{e}}_j$。再以共享比较网络$G$形成局部差异表示，最后分类：
\begin{equation}
 p(r\mid e,c)=\operatorname{softmax}
 H\left[
 \sum_iG(\bm{e}_i,\widetilde{\bm{c}}_i);
 \sum_jG(\bm{c}_j,\widetilde{\bm{e}}_j)
 \right].
 \label{eq:nlp-dialogue-decomposable}
\end{equation}
$F,G,H$的参数由三分类交叉熵更新。基础模型主要对词表示做匹配，词序和复杂范围关系的表示有局限；可增加句内注意力，或用现代联合编码器替换分类部件，但输入证据协议与标签含义仍需保持。

取主张“当前读书会在4月18日举行”。证据甲为共同通知的改期句，应支持；证据乙为另一份明确假设的权威更新“当前活动改为4月19日”，在该更新生效的独立测试情形下应反驳；证据丙为“海棠楼201正在检修”，对活动日期不提供确定信息。假设模型输出三类概率分别为$(0.90,0.04,0.06)$、$(0.05,0.90,0.05)$、$(0.15,0.10,0.75)$，按最大值可恢复三种关系；这些概率是教学输出，不是规则证明。

若证据分成“ML-0418由甲社团举办”和“甲社团的该次读书会改至4月18日”，单句可能都不足以完成对象绑定，联合判断才有机会恢复支持。检索器没有取得第二句和分类器看见两句却未能组合，应计为不同错误。在测试中还可提供正确证据测分类能力，再从全资料运行测端到端能力；不能用前者代替后者。

\subsubsection{规则与结构化事实核验}
\label{sec:nlp-dialogue-structured-check}

日期、数值和明确字段可以先规范化再比对。规则输入是带对象标识、属性、时间和单位的主张记录，以及通过只读查询取得的证据记录；规则来自业务定义与数据模式。只有主体、属性和适用范围一致，才比较值。相同则输出匹配，不相容则输出冲突，缺字段或范围未对齐则输出无法判定，并附上原因。

例如主张“报名可以在4月16日18:00提交”经时间规范化得到$t=2025\text{-}04\text{-}16\ 18{:}00$。证据条件为$t<t_{\mathrm{deadline}}$，代入边界时刻得到假，因此与通知冲突。主张“到场人数为0.12千人”和证据“120人”规范到相同单位后匹配；“120组”则因缺少每组人数而无法判定，不能仅比较数字120。活动原计划4月17日、当前4月18日属于不同时间属性，保留属性后不冲突。

这类规则能精确覆盖已编码的比较，却不能检查未列出的隐含因果。例如“人数增长是因为改到了周五”增加了因果主张，通知和人数表都不足以支持。规则应返回不在支持范围内，交给有相应证据的后续判断，而非因日期与人数分别正确就判整句为真。学习式关系模型可以处理更丰富的语言，但也应接受同一组单位、边界时刻和对象反例的测试。

\subsection{核验结果与回答边界}
\label{sec:nlp-dialogue-answer-boundary}

核验结束后，需要把逐项关系转成可读回答。对日期、地点和人数三个主张，整句弃答会丢掉已获支持的两项；局部弃答则可以恢复“通知说明活动于4月18日在海棠楼201举行；通知未给出实际到场人数”。若已经获得独立签到统计S，可修订为“签到统计记录为120人”，并引用S而非通知。修订器以原回答、主张清单、证据与允许修改范围为输入，可以用规则替换或条件编辑生成；训练时需要与证据一致的修订目标，提示式使用时也应保留相同约束。

是否接受模型判定，还需在独立开发样本上选择阈值。设四条教学主张的支持置信为$0.95,0.85,0.75,0.40$，独立核验结果依次为对、错、对、错。阈值$0.8$接受两条，覆盖率为$2/4$，接受样本错误率为$1/2$；阈值$0.9$只接受一条，覆盖率为$1/4$，接受样本错误率为0。后一项只描述这四个样本，不能当作上线零错误保证。阈值选择和最终测试必须分离，并同时报告覆盖与风险，而不只报告保留下来的高分回答。

错误原因还应区分材料本身未提供、检索未找到、模型未识别和来源相互冲突。来源冲突时保留版本和适用时点，无法确定哪个生效就应陈述冲突或请求补充；不能把模型偏好的来源称作已核实的新版本。修订后重新拆解并检查是否引入新主张，例如把“无报名记录”扩写为“您没有报名”就越过了证据范围。自我置信表达或第二次无依据生成不代替独立样本校准与人工复核。

\section{任务型对话}
\label{sec:nlp-dialogue-task}

若用户只是问活动何时举行，查通知即可；若要核验自己一行人的报名情况，答案还取决于身份、查询条件和业务记录。任务型对话把这些条件保存在跨轮需求状态中，再据此选择行为。理解负责本轮说了什么，状态负责目前哪些要求有效，策略负责下一步做什么，回复负责如实解释结果。单轮回答自然不表示这四项都正确。

本节固定如下教学业务，所列记录均不是共同通知本身提供的事实。咨询时点为2025年4月17日，用户已经通过登录认证，身份令牌为U17，目标是核对已有报名，不是新增报名。只读接口\code{check\_registration}接受活动编号、认证身份、活动日期和同行人数。假设其快照中有一条U17记录：活动ML-0418，日期4月18日，人数2，提交时点4月16日17:30，状态有效。它早于严格截止时点，允许回答“查到按时提交的两人报名记录”，但不允许声称系统刚办理了新报名或保证所有入场条件均已满足。

三轮输入固定为“帮我核对读书会的报名记录”“4月17日，两人”“说错了，是4月18日，人数不变”。第一轮缺日期与人数，第二轮条件下查无匹配，第三轮更正日期后查到已有记录。图\ref{fig:nlp-dialogue-state}还画出接口超时的独立分支；超时不能解释成查无记录。主线采用JointBERT式轮次理解、TRADE式整状态预测与规则行为选择的教学组合。原TRADE直接读取对话历史，不以JointBERT的槽位标签作为必需输入；本例并列记录局部理解与整状态，便于找到两者不一致的环节。

\begin{figure*}[htbp]
  \centering
  % Book-styled flowchart; solid arrows are data, dashed arrows update state.
\begin{tikzpicture}[
  every node/.style={font=\figurefont\small,align=center},
  utterance/.style={book node,text width=38mm,minimum height=16mm,
    fill=FigureInput!10,draw=FigureInput},
  state/.style={book node,text width=42mm,minimum height=16mm,
    fill=FigureModel!10,draw=FigureModel},
  result/.style={book node,text width=48mm,minimum height=16mm,
    fill=FigureOutput!10,draw=FigureOutput},
  flow/.style={book arrow,rounded corners=1.2mm},
  update/.style={flow,dashed,draw=FigureModel},
  labeltext/.style={text=BookMuted,font=\figurefont\small}
]
  \node[labeltext] at (0,14mm) {用户输入};
  \node[labeltext] at (48mm,14mm) {当前需求状态};
  \node[labeltext] at (104mm,14mm) {行为与观察};
  \node[utterance] (u1) at (0,0) {第1轮：核对报名\\日期、人数未给出};
  \node[state] (s1) at (48mm,0) {日期：缺失\\人数：缺失};
  \node[result] (a1) at (104mm,0) {缺必需字段\\询问日期与人数};
  \node[utterance] (u2) at (0,-28mm) {第2轮：新增约束\\4月17日、两人};
  \node[state] (s2) at (48mm,-28mm) {日期：4月17日\\人数：2};
  \node[result] (a2) at (104mm,-28mm) {只读查询：查无匹配\\核对是否指4月18日};
  \node[utterance] (u3) at (0,-56mm) {第3轮：更正日期\\4月18日、人数不变};
  \node[state] (s3) at (48mm,-56mm) {日期：4月18日\\人数：2（保持）};
  \node[result] (a3) at (104mm,-56mm) {只读查询：匹配1条\\核验两人已按时提交};
  \foreach \i in {1,2,3} {
    \draw[update] (u\i.east) -- (s\i.west);
    \draw[flow] (s\i.east) -- (a\i.west);
  }
  \draw[flow] (a1.south) -- (104mm,-14mm)
    -- (0,-14mm) -- (u2.north);
  \draw[flow] (a2.south) -- (104mm,-42mm)
    -- (0,-42mm) -- (u3.north);
  \node[labeltext,fill=white,inner sep=1mm] at (48mm,-14mm) {回复：补全必需条件};
  \node[labeltext,fill=white,inner sep=1mm] at (48mm,-42mm) {回问：让用户确认修改};
  \node[book node,text width=110mm,minimum height=12mm,
    fill=BookAmber!10,draw=BookAmber] (failure) at (53mm,-84mm)
    {独立故障分支：TIMEOUT不是查无记录\\说明失败，按预算重试或回问是否稍后继续};
  \draw[flow,draw=BookAmber] (a2.east) -- (138mm,-28mm)
    -- (138mm,-84mm) -- (failure.east);
  \node[labeltext] at (53mm,-98mm)
    {固定上下文：认证U17、活动ML-0418；不新增或修改报名};
\end{tikzpicture}

  \caption{只读报名核验的三轮对话。用户目标为核对U17在ML-0418的已有报名；实线传递话语、动作和接口结果，虚线表示状态写入或更正。第二轮查无匹配后回问，第三轮只覆盖日期而保持人数。接口故障另走重试或说明失败的分支，不产生报名结论。}
  \label{fig:nlp-dialogue-state}
\end{figure*}
\FloatBarrier

\subsection{轮次理解}
\label{sec:nlp-dialogue-turn}

轮次理解的输出可以包括意图、槽位、对话行为和回应对象。意图如“核对报名”，槽位如日期和人数，行为如提供条件或更正，回应对象则说明“不是明天，是下周”否定的是哪次提议。本节文本输入保留说话者、轮号和此前系统问题，日期解析以已给定的咨询时点为参照。声音中的停顿与重音可能有帮助，但那属于第\ref{chap:nlp-speech-tasks}章的额外输入条件。

\subsubsection{意图分类与CRF槽位标注}
\label{sec:nlp-dialogue-intent-crf}

流水线先预测句级意图，再给词元序列赋BIO标签。训练样本包含话语、意图类别、槽位字符区间，以及分词到字符的映射；若分类器或标注器使用历史，需要在训练和测试中固定相同历史范围。意图模型用分类交叉熵，CRF用正确标签路径的负对数似然，标签转移和输入特征的得分共同决定合法路径。分类与CRF计算分别见第\ref{chap:nlp-classification-matching-organization}、\ref{chap:nlp-text-analysis}章，此处增加的是将两个输出恢复成一次交流记录。

把第二轮按“4月／17日／，／两／人”分词，正确标签为\code{B-DATE I-DATE O B-PARTY I-PARTY}。CRF解码后把相邻日期标签合回原文“4月17日”，把人数标签合回“两人”，再规范化为日期2025-04-17和正整数2。假设意图分类在“核对报名、取消报名、查询通知”上给出$(0.8,0.1,0.1)$，选核对报名；假设正确标签路径分数5，错误路径分数3，在仅有这两条路径的教学空间中，正确路径概率为$e^5/(e^5+e^3)=0.881$。实际CRF的分母包含全部合法路径，不只是最佳两个候选。

若意图误判为取消报名，而实现又据意图限制允许槽位或动作，日期即使抽对也可能被送入错误流程。因此本节只读服务在业务层禁用取消动作，不把分类结果当权限。第三轮同时出现旧值否定与新值时，普通DATE标签只能识别日期位置，不能表示哪个值应撤销。更正、撤回和被更正槽位需要另行标注或显式规则；将两个日期都抽出不是状态更新算法。

\subsubsection{JointBERT联合意图与槽位}
\label{sec:nlp-dialogue-jointbert}

JointBERT让意图与槽位共享BERT表示。句首\code{[CLS]}的表示$\bm{h}_0$进入意图softmax，各原词的首个子词表示进入槽位softmax：
\begin{equation}
\begin{aligned}
 p(c\mid u)&=\operatorname{softmax}(\bm{W}_{c}\bm{h}_0+\bm{b}_{c}),\\
 p(s_i\mid u)&=\operatorname{softmax}(\bm{W}_{s}\bm{h}_{a(i)}+\bm{b}_{s}),\\
 L&=-\log p(c^*\mid u)-\sum_i\log p(s_i^*\mid u).
\end{aligned}
\label{eq:nlp-dialogue-jointbert}
\end{equation}
其中$a(i)$是原词$i$的首子词位置。其余子词不重复承担该原词的监督；预测后依据分词映射恢复原词标签和原文跨度。原文也比较了在槽位输出上增加CRF的版本，但基础模型的槽位预测是逐位置分类，不能把联合意图训练误说成已经有标签转移约束（见文献说明）。

固定上一小节第二轮的意图与槽位标注，设正确意图概率0.8，五个正确BIO标签概率均为0.9，则示例总损失为$-\log0.8-5\log0.9=0.750$。两个头的梯度都更新共享编码器，这是与分别训练两个模型的直接差别；它不保证两个任务永远互相促进。无CRF时若出现句首\code{I-DATE}，恢复器应按预先规定的合法化规则处理或拒绝，而不能在评测时临时选最有利的解释。

第一轮可以预测核对意图而无日期、人数槽位；第二轮提供两个槽位；第三轮日期槽位为“4月18日”，人数没有新值。“人数不变”对应保持指令，需在本书教学组合中增加\code{KEEP(party)}行为解析，不能把未再次抽到人数评为人数已消失。多意图、否定范围、回应对象和更正行为都不在单意图加普通BIO标签的默认输出中；扩展它们需要新的目标和样本。评价时分别检查本轮标签和跨轮状态，防止把局部高F1误当完整理解。

\subsubsection{模式约束的语义记录生成}
\label{sec:nlp-dialogue-semantic-record}

另一种做法是直接生成结构记录。给模型当前轮、限定历史、字段模式和示例，规定\code{intent}、\code{operation}、\code{slot}、\code{old\_value}、\code{new\_value}及\code{source\_turn}的类型；训练时用人工标注记录的词元负对数似然，提示式实现则固定示例和模型版本。语法限制可以在解码时屏蔽非法键、枚举值或括号，生成后再用结构解析器检查。第\ref{chap:nlp-text-generation}章的受约束生成负责格式，本节负责字段的交流含义。

第三轮应恢复两个操作：日期\code{REPLACE(2025-04-17,2025-04-18)}，人数\code{KEEP}。新日期的原文依据在第3轮，旧日期的依据在第2轮；\code{old\_value}不能凭空猜测。若上一轮同时讨论活动日期与报名截止，“改成周五”未说明修改对象，就应输出待澄清，而不是覆盖两个字段。未知值用显式状态表示，不能生成字符串\code{"unknown"}后当成合法日期去查询。

开放值生成避免了为所有名称预列类别，但会增加格式合法、含义错误的记录。例如\code{REPLACE(date,2025-04-18)}键值均合法，若原话是否定“不是4月18日”，理解仍然相反。恢复时逐一核对字段类型、来源跨度、否定范围和历史绑定。与标签路线比较应使用同一轮次及同一语义参考，分别报告槽位识别、操作识别和完整记录匹配；额外生成字段属于错误，不因结构丰富而加分。

\subsection{需求状态的更新}
\label{sec:nlp-dialogue-state}

需求状态保存当前有效的领域—槽位—值，而不是历史上出现过的所有词。MultiWOZ提供多领域对话与状态标注，SGD进一步用服务和槽位描述约束接口，并区分训练中见过和未见过的服务；它们提供输入与监督协议，不意味着换一个服务名称就能无条件迁移（见文献说明）。

本书把四种情况分开：尚未提及为缺失，明确说不知道为未知，表示任意值都可接受为不限定，明确拒绝某个值为否定约束。当前活动日期、同行人数属于会话状态；已确认且允许长期保存的偏好才可能进入持久记忆；训练好的权重则是模型参数。这三个对象不能都称为“记住了用户”。

\subsubsection{规则槽位写入与撤回}
\label{sec:nlp-dialogue-state-rule}

规则更新器接收旧状态$B_{t-1}$和本轮已解析操作$\Delta_t$，产生$B_t=U(B_{t-1},\Delta_t)$，同时保留来源轮号。规则来自固定字段模式与业务含义，不需要训练状态预测参数。表\ref{tab:nlp-dialogue-state-update}给出本节使用的转换；冲突不采用最后一个词自动覆盖，而是留下待澄清标记。

\begin{table}[htbp]
\centering
\begin{tabular}{p{22mm}p{74mm}}
\toprule
操作 & 状态转换\\
\midrule
新增 & 旧槽位缺失时写入值与来源；已有不相容值则待澄清\\
更正 & 明确指向旧槽位时覆盖，并保留旧值为历史而非有效值\\
保持 & 不改变原值及其来源；本轮可另记保持指令\\
撤回 & 删除有效正值；若只拒绝某值，则写否定集合而非不限定\\
未知／不限定 & 分别写入不同标记；不可按同一种空字符串处理\\
\bottomrule
\end{tabular}
\caption{需求状态的规则更新。缺失是没有当前值，撤回是有来源的交流事件，两者在最终正值表中可能相同，在历史记录中仍不同。}
\label{tab:nlp-dialogue-state-update}
\end{table}

活动编号和登录身份由已确认上下文提供，以下只演算用户可修改的两个槽位。初始$B_0=\varnothing$；第一轮$B_1=\varnothing$；第二轮$B_2=\{(\text{日期},4/17),(\text{人数},2)\}$；第三轮$B_3=\{(\text{日期},4/18),(\text{人数},2)\}$。两项$B_2$值来自第2轮，$B_3$日期来自第3轮，人数仍来自第2轮。第三轮变化集只有一个覆盖操作，不能因本轮没有新人数而把2清空。

增量规则容易审计每次改动，但前面误写的值可能长期保留；从历史重构完整状态有机会改正旧错，也可能无故遗漏未重提的约束。可在每轮同时核对完整状态和变化项。若用户另说“不查了”，应更新目标与流程状态并停止查询；这不是把日期改成空值后继续追问日期。规则的可靠性取决于输入操作是否正确，不能弥补未识别的否定和指代。

\subsubsection{TRADE复制式状态生成}
\label{sec:nlp-dialogue-trade}

TRADE以选定范围的用户话语和系统回复为输入，用双向GRU编码历史。对每个预先给定的领域—槽位对，以两者嵌入之和启动共享GRU解码器，独立生成该槽位的值；第一个解码步的历史注意摘要还进入三分类门控。门控为\code{none}时不输出有效值，为\code{dontcare}时输出用户不限定，只有\code{ptr}时采用生成的值。这里\code{ptr}不意味着只能复制，值解码同时混合词表分布和历史复制分布\scite{nlp-dialogue-wu2019}

令历史位置表示为$\bm{h}_i$，当前解码状态为$\bm{g}_{j,k}$，位置注意力$\alpha_{j,k,i}$由内积softmax得到。相同词在不同历史位置出现时，其复制概率相加：
\begin{equation}
 p(y\mid j,k)=
 \gamma_{j,k}p_{\mathrm{vocab}}(y\mid j,k)
 + (1-\gamma_{j,k})\sum_{i:x_i=y}\alpha_{j,k,i}.
 \label{eq:nlp-dialogue-trade}
\end{equation}
$\gamma_{j,k}$由当前解码状态、输入词嵌入和历史摘要经sigmoid计算。门控交叉熵与正确状态值词元的负对数似然加权训练共享参数，监督是每轮完整状态，不只是本轮新槽位。解码到结束符后恢复值，再做日期和整数规范化；规范化是本节服务适配，不应冒称原模型天然输出本书的日期格式。

第三轮解码日期时，假设$\gamma=0.2$，词表对“4月18日”的概率0.1，历史对应位置的复制概率0.7，则混合值为$0.2\times0.1+0.8\times0.7=0.58$。旧“4月17日”若词表概率0.2、复制概率0.1，则为0.12，解码选择新日期。人数槽位仍可从第2轮复制“两人”，经规范化保留2。新出现的名称即使不在固定词表中，也可能从输入位置取得概率；但若门控先误判\code{none}，复制分数再高也不会写入状态。

原TRADE的三个门控标签没有逐一表达本书的未知、撤回事件和否定集合。撤回后无有效正值可以映射到\code{none}，但“用户从未提及”和“用户明确撤销”仍需额外历史；未知不能无说明地映射成\code{dontcare}。同样，复制旧轮次中的日期会生成完全合法却已失效的值。未见领域测试仍需提供领域—槽位输入，且原文的迁移动机依赖已有相似槽位；不能把复制新名称推广为对任意未见槽位含义的理解。

\subsubsection{SimpleTOD完整状态序列生成}
\label{sec:nlp-dialogue-simpletod}

SimpleTOD将对话历史$C$、需求状态$B$、数据库摘要$D$、系统行为$A$和去词汇化回复$S$按顺序放入带边界标记的序列。用预训练因果语言模型初始化，以各位置的下一个词元负对数似然训练，状态、行为和回复的监督由任务数据提供。边界标记使解析器能区别“用户曾说过的日期”和“当前状态字段”，但标记本身不保证后者正确（见文献说明）。

使用时先由$C$生成$B$并在状态结束符停下，解析、校验后实际查数据库取得$D$，再以$[C;B;D]$生成$A,S$。原文的$D$包含匹配条目数量以及相应实验设置下的预订状态信息；它是查询返回的摘要，不是让模型随意续写后信任的“业务结果”。具体返回行保留在程序侧，用于最后槽位回填。

第三轮教学序列可读作“历史：三轮；状态：日期4/18、人数2；数据库：匹配1条；行为：告知已有有效记录；回复：查到您在[date]的[party]人报名记录”。其中“匹配1条”是一次调用后的观察，不是参加人数1。若语言模型遗漏人数，类型校验可能仍通过，但查询前的必需槽位检查必须拦截；若把日期续写为4/17，则会再次得到空结果，不能靠流畅回复覆盖这个错误。

与规则增量更新和TRADE相比，完整序列生成共享更多任务参数，也让状态、动作和回复错误相互影响。三轮比较需固定历史、服务快照和状态规范化，并记录数据版本及标注修订；不能拿一种模型的修正标签与另一种模型的原标签直接比较。状态的联合目标准确率要求每轮全部槽位正确，变化项指标只检查本轮应改动部分；两者都不能取代实际调用和最终回复验收。

\subsection{响应选择与业务调用}
\label{sec:nlp-dialogue-action}

状态完整后仍需决定是否以及怎样调用服务。业务模式给出字段类型、必需条件、可用动作和停止条件，策略则从询问、澄清、查询、确认等行为中选择当前一步。复杂环境中的长程计划由第\ref{chap:nlp-solving-and-execution}章展开；这里先把一次服务调用的前后边界说明白。

\subsubsection{有限状态与槽位驱动策略}
\label{sec:nlp-dialogue-rule-policy}

本节策略由服务定义者写定：缺必需字段则询问；字段冲突或指代不明则澄清；参数齐备且认证有效才允许只读查询；成功取得与目标相符的结果后告知并停止。允许动作掩码将不可用动作的分数设为$-\infty$，但掩码只约束选择空间，不判断哪个允许动作最有帮助。由于已过报名截止且服务只读，本例任何轮次都不允许生成新增报名调用。

第一轮日期与人数均缺失，动作是“请问您核对哪天、几人的报名记录”。第二轮将两个槽位与可信上下文回填为\code{(ML-0418,U17,2025-04-17,2)}，执行只读查询。接口返回\code{OK}且记录列表为空，规则转到条件核对：“按4月17日、两人未查到匹配记录；通知中的活动日期为4月18日，您要核对这个日期吗？”不能默默把用户日期改成通知日期再声称其原请求已有结果。

第三轮覆盖日期、保持人数，参数变成\code{(ML-0418,U17,2025-04-18,2)}，查询返回一条有效记录。规则另验提交时点$17{:}30<18{:}00$，形成告知动作并结束本次核验。若返回\code{TIMEOUT}，则没有可用的记录判断，允许动作是说明失败、按预算重试或请求稍后继续；若用户改为“只想看通知”，应改变目标并取消不再必要的身份查询。流程走到查询节点不是任务成功，须等实际观察满足目标。

\subsubsection{监督动作预测与结构化行动生成}
\label{sec:nlp-dialogue-learned-policy}

规则多时，可以从带行为标注的对话学习策略。分类路线将历史摘要、当前状态和此前接口观察编码为$\bm{v}_t$，以$\operatorname{softmax}(\bm{W}\bm{v}_t)$预测离散动作，最小化正确动作交叉熵；参数仍由类型化槽位回填。生成路线以相同输入预测动作名和参数组成的结构记录，以记录词元负对数似然训练，并限制动作名、字段及枚举值。两条路线都需要与实际服务模式一致的标注。

第一轮若未加掩码时$\Pr(\text{询问})=0.7,\Pr(\text{查询})=0.3$，缺日期和人数使查询不可用，恢复的动作只能是询问。第二轮完整状态上若分类器选择查询，回填同一组四元参数；生成器也必须解析为相同参数，才能和规则策略公平比较。若生成器写出人数\code{"两"}而模式要求整数，先由声明的规范化器变成2或拒绝，不能依赖服务容忍所有格式。

未见服务需要把名称、描述、必需字段和前置条件提供给模型，并单独评价字段映射；曾学会\code{date}不表示知道新接口的日期是出发日还是报名日。强化学习可用完整任务结果和交互成本优化策略，但需要固定环境、奖励与探索权限，不能由几个监督动作样本推导其性能。无论采用何种训练，行动分数不是业务成功概率，生成合法调用文本也不是已经执行；调用日志和返回结果必须另外保存。

\subsection{结果解释与交流恢复}
\label{sec:nlp-dialogue-response}

需求状态描述用户当前要什么，业务状态描述环境实际有什么，回复必须在两者之间建立有依据的对应。用户更正日期只改变需求状态，不会自动修改报名记录；数据库更新也不等于用户目标随之改变。$\tau^2$-bench进一步考察用户和智能体均可用工具改变共享环境的条件，提醒评测不能把用户当成只提供静态文本的角色；本节只读快照是比它更受限的教学设置（见文献说明）。

\subsubsection{去词汇化模板与槽位回填}
\label{sec:nlp-dialogue-template}

模板把结果类型与语言表达分开。设计者预先定义成功、空结果和故障三类模板，指定每个占位字段来自接口返回还是当前查询条件，日期显示规则与人数单位也固定。成功模板可为“查到[date]的[party]人有效报名，提交于[submitted]”；空结果模板为“按[date]、[party]人未查到匹配记录，请核对条件”；故障模板不包含任何记录是否存在的断言。

第二轮的\code{OK,records=[]}选择空结果模板，日期和人数来自本次请求，不能填快照中另一条记录。第三轮取得匹配行后，成功模板填“4月18日”“2”“4月16日17:30”，再依据截止规则补充“该提交时点早于通知规定的截止时点”。回填前检查返回行的身份与活动绑定、请求和结果的日期一致、所有必需占位符均已消除。若用户已经改目标而旧异步结果才到达，应按请求版本丢弃或标注过期，不能回复给新目标。

这三轮中第一轮没有事实答案，仍完成了必要询问；第二轮没有匹配记录，仍可正确完成一次查询并帮助用户恢复；第三轮才达到已有报名核验目标。模板没有训练损失，质量依赖结果映射和人工覆盖测试。它限制表达变化，却不能阻止上游错状态；只检查没有残留占位符也不能发现“未查到”被错误改成“未报名”。

\subsubsection{SC-LSTM语义条件回复}
\label{sec:nlp-dialogue-sclstm}

语义控制LSTM（semantically controlled LSTM, SC-LSTM）从对话行为和槽位语义生成回复。输入控制向量$\bm{d}_0$由行为类型及槽位值特征组成，文本训练目标经过去词汇化；模型学习哪些表达对应哪些语义，而不要求逐词人工对齐。它在普通LSTM外增加读取门，让尚待表达的语义随生成逐渐衰减：
\begin{equation}
\begin{aligned}
 \bm{r}_i&=\sigma(\bm{W}_{wr}\bm{w}_i+
                   \alpha\bm{W}_{hr}\bm{h}_{i-1}),\\
 \bm{d}_i&=\bm{r}_i\odot\bm{d}_{i-1},\\
 \bm{c}_i&=\bm{f}_i\odot\bm{c}_{i-1}
       +\bm{\imath}_i\odot\widetilde{\bm{c}}_i
       +\tanh(\bm{W}_{dc}\bm{d}_i).
\end{aligned}
\label{eq:nlp-dialogue-sclstm}
\end{equation}
其中$\bm{\imath}_i$为输入门，$\bm{f}_i$为遗忘门；普通循环网络机制见模型卷。语义向量进入细胞状态，读取门则学习哪些信息可以从待表达条件中移除\scite{nlp-dialogue-wen2015}

训练以参考词元的负对数似然为主体，并惩罚结束时残留的语义以及过于突变的门控，避免未表达完就停止或一次关闭过多特征。原方法先从前向分布随机采样多个候选，再结合前向、反向模型分数和槽位错误惩罚重排。这里不是在解码前硬性规定每个槽位只能出现一次，仍须检查候选是否遗漏或重复。

对第三轮结果，设三个待表达语义为日期、人数、提交时点，初始简化向量$(1,1,1)$。假设输出日期附近的读取门为$(0.1,1,1)$，更新成$(0.1,1,1)$；表达人数时门为$(1,0.1,1)$，更新成$(0.1,0.1,1)$。这是软保留量的教学演算，不是概率为零一的已完成证明。若候选只生成“查到[date]的[party]人报名”，缺少提交时点，三个必需槽位的遗漏率为$1/3$；补齐候选才可参加最后选择。

最终仍用真实记录回填日期、人数和时点，并复核“有效”“按时”等关系是否成立。原文的精确槽位匹配对可去词汇化字段有效，对二元属性和“不限定”这类可用多种说法表达的条件并不充分。因此语义控制提高可学习的表达能力，却没有取消外部检查。接口失败、用户更正和目标变化先更新输入行为与状态，不能仅让语言模型换一种安慰语气；完整任务成功也不能由某一轮回复偏好替代。

\section{开放对话}
\label{sec:nlp-dialogue-open}

不是每次交流都有一张待填写的业务表。用户可能希望延续一个知识话题、聊自己的兴趣，或说出一次不顺利的经历。开放对话的目标较难用单一接口结果表示，但仍需区分知识、人物资料和当下情绪的用途。知识为可查事实提供依据，人物资料约束某个主体的经历与偏好，情绪线索帮助判断回应方式；三者不能互相代替。

多方会话还需保留说话者和回应对象。“我第一次参加”属于当前发言人，不一定属于群聊中上一位同名成员；某人说“我也是”可能回应经历，也可能只回应感受。下面各方法均在给定历史范围和身份绑定下工作，评价时将事实依据、跨轮连贯与交流用途分开。

\subsection{知识支持的交流}
\label{sec:nlp-dialogue-knowledge}

知识支持对话需要在历史中找到尚待回答的属性。用户先问活动时间，下一轮问“那地点呢”，第二轮应沿用同一活动对象而切换待答属性，不能重复完整的时间回答。Wizard of Wikipedia把人类交流中选用的知识句与回复关联起来，使资料取得、细粒度知识选择和回复形成可以分别检查。以下用共同通知构造微型例子，方法来源见本章末文献说明。

\subsubsection{检索式Transformer Memory Network}
\label{sec:nlp-dialogue-retrieval-memory}

检索式Transformer Memory Network先用信息检索取得知识句，再据知识和历史给候选回复排序。原文分别以初始话题及最近两轮为查询，检索百科条目的首段，将候选拆成附带标题的句子。此处可用第\ref{sec:nlp-dialogue-dpr}节的词汇或稠密检索替换召回器，但必须注明替换，不能把对整部百科做一次神经注意力当作原输入过程。

用共享Transformer编码历史为$\bm{q}$、知识句为$\bm{m}_j$，得到
\begin{equation}
 \alpha_j=\frac{\exp(\bm{q}^{\top}\bm{m}_j)}
                  {\sum_k\exp(\bm{q}^{\top}\bm{m}_k)},
 \qquad
 \bm{v}=\bm{q}+\sum_j\alpha_j\bm{m}_j.
 \label{eq:nlp-dialogue-retrieval-memory}
\end{equation}
另一Transformer编码回复候选为$\bm{r}_l$，按$\cos(\bm{v},\bm{r}_l)$取最高者。原检索式模型的候选来自训练集回复，训练以正确回复为正例、批内其他回复为负例做交叉熵；细粒度知识权重可在回复目标下学习，不能因此声称每条最高权重知识都有人类监督。候选限制也意味着它不能凭空输出候选集中不存在的新事实表述。

取知识$K_1$为“活动4月18日14:00举行”，$K_2$为“地点海棠楼201”。教学向量规定$\bm{m}_1=(1,0),\bm{m}_2=(0,1)$。第一轮问时间，令$\bm{q}=(1,0)$，则知识权重为$(0.731,0.269)$，$\bm{v}=(1.731,0.269)$。时间回复和地点回复的向量分别取$(1,0),(0,1)$，余弦分数约为0.988与0.154，选择时间回复。下一轮“那地点呢”若被正确结合历史编码为$(0,1)$，权重和两个回复得分对调，恢复地点回复而不重复时间。

例子中的向量是规定值，只演示融合与排序。真实模型可能召回错活动、选错属性，或知识选对而候选回复只说“那里很好”；应分别计算知识召回、知识选择和回复选择。最高注意权重便于检查模型偏向，却不是逐句引用证明。来源冲突和版本更新仍回到第\ref{sec:nlp-dialogue-verification}节处理，不能为了维持话题而沿用失效知识。

\subsubsection{生成式Transformer Memory Network}
\label{sec:nlp-dialogue-generative-memory}

生成式版本允许逐词形成新回复，但知识如何进入解码器必须说清。两阶段路线单独训练知识选择器，选定最高分知识句，再把该句与历史连接交给另一个编码器，训练生成器预测参考回复。知识选择用人类所选句子的交叉熵，回复用词元负对数似然；两个模型可分别改进，也会发生前一阶段错误向后传播。

原文的端到端版本共享历史和候选知识的Transformer编码器，将候选句表示汇总后计算知识注意分布，但随后仍取最高分的一句，把它的完整词元表示与历史表示连接供解码器访问。它不是式\eqref{eq:nlp-dialogue-retrieval-memory}中所有句向量的软加权生成，也不是FiD同时访问全部片段。训练目标可加辅助知识监督：
\begin{equation}
 L=(1-\lambda)L_{\mathrm{response}}+\lambda L_{\mathrm{knowledge}}.
 \label{eq:nlp-dialogue-generative-memory}
\end{equation}
硬选择不会仅因模型被称作端到端就变成可微软混合；辅助损失直接给知识注意提供监督，解码损失则训练实际采用的表示和生成过程。

固定上例第二轮，设选择器给地点句概率0.8，人类知识标签为地点句，则知识损失$-\log0.8=0.223$。生成器采用该句和含时间回答的历史，输出“地点是海棠楼201”，而不是从固定回复候选中选择。原论文用束宽5的搜索；教学恢复可同样保留五条候选，逐条检查地点、活动绑定和未获支持的追加语。若候选为“海棠楼201，现场免费提供资料”，前半句有依据，后半句没有，不能因知识选择正确就整体通过。

两阶段的重新编码允许选定知识与历史在编码时交互，端到端版本则复用独立编码结果；两者都要处理只选一句时的知识覆盖不足。需要组合多片段可比较RAG或FiD扩展，并明确已改变模型接口。跨轮评价还要检查省略对象的解析、用户已知内容是否被无谓重复以及话题切换是否被尊重，不能把每一轮当成互不相关的事实问答。

\subsection{社交与个性化交流}
\label{sec:nlp-dialogue-personalization}

人物条件使同一句问题在不同主体下需要不同回答。Persona-Chat给交流双方指定人物资料，比较带资料与不带资料的回复；这种设定评价的是与指定人物的一致性，并不证明虚构人物在现实中存在。用于真实用户时，还要规定哪些资料允许使用、从何而来以及何时失效。

\subsubsection{Profile Memory Network人物条件回复}
\label{sec:nlp-dialogue-profile-memory}

Profile Memory Network把人物条目作为独立记忆，不把它们与对话历史无区别地合成一句长文本。排序式版本用嵌入汇总表示历史$\bm{q}$和每条人物资料$\bm{p}_j$，由相似度softmax得到权重，形成$\bm{q}^{+}=\bm{q}+\sum_j\alpha_j\bm{p}_j$，再按与候选回复的相似度排序。它沿用StarSpace式正确回复与负回复的排序训练；相较直接连接人物资料的基线，区别在于可以按当前话题选择资料条目（见文献说明）。

生成式版本用LSTM编码历史，最后状态初始化解码器；人物条目则按词频加权汇总为独立向量，形成矩阵$\bm{P}$。解码时依当前隐状态计算$\operatorname{softmax}(\bm{P}\bm{W}_a\bm{h}_i)$，所得人物摘要与词元输入结合，预测下一词。训练使用人物资料、历史与参考回复的词元负对数似然。把全部资料直接前置到序列的普通序列到序列模型是可比较基线，但不是上述逐步人物注意机制。

教学角色甲的资料为“喜欢机器学习”“去年做过读书会志愿者”，角色乙为“喜欢机器学习”“这次第一次参加读书会”。问题“以前来过吗”下，两个角色都可能对兴趣条目产生较低权重，经历条目才决定回答。若甲两条权重为$(0.2,0.8)$，可以恢复“我去年做过志愿者”；乙则应恢复“这是我第一次参加”。生成甲的经历前必须先确定正在扮演甲，不能把甲乙全部资料一起做注意力，再从中选一个较相关的经历。

排序输出的完整候选和生成输出的词元都应对照人物条目检查。资料只说做过志愿者，不支持“我去年主持了活动”；对方说“原来你经验丰富”也不能自动变成人物记忆中的已确认专业资历。对人物一致性、话题相关性和现实事实真实性应分开评分，资料缺失时允许不作具体经历声明。

\subsubsection{用户记忆检索与规则更新}
\label{sec:nlp-dialogue-user-memory}

真实个性化还需要一个持久记录协议。这里采用本书教学基线：每条记忆保存主体标识、属性、值、来源轮次、确认时点、有效范围和使用许可；明确陈述才进入候选写入，推测保留为未确认而不直接持久化。语义编码器提供相关性分数，其参数可来自第\ref{chap:nlp-classification-matching-organization}章的匹配训练；身份、有效期和更新优先级由规则决定。存储实现与访问控制属于系统卷，本节只规定交流语义。

假设两位同名“林宁”的身份分别为U17、U28。记录M1为U17在4月10日确认“默认先给简短回复”，M2为U28偏好“完整推导”，M3为U17在4月17日说“这次请详细解释”。当前请求来自U17。即使M2语义相似度0.95高于M1的0.80，也因主体不同被过滤；M3只对本次任务有效，优先于长期默认M1，但不删除M1。回复时采用详细解释，任务结束后下一次请求仍可用M1。

若U17明确说“以后都改成详细解释”，规则在确认后把M1标记失效，写入有新来源的长期记录；若只说“这个回答太短”，尚不能推出所有未来任务都需详细。发生同时间相互矛盾的明确指令时，保留冲突并询问，不以向量相似度裁决。多方会话中还需检查发言人与被描述主体，例如U28说“林宁想要简短版”不能未经绑定就覆盖U17的确认记录。

恢复给生成器的不是一包相似文本，而是已筛选的有效条件及来源。评价可在同名、过期、更正和临时覆盖样本上分别检查命中、误用、错误写入与回复用途。用户偏好改变不是模型参数在线更新；模型自己在上一轮说过“您可能喜欢详细解释”，也不是一条用户证据。

\subsection{共情与支持性交流}
\label{sec:nlp-dialogue-empathy}

“我没赶上报名”可能表达失落，也可能只是请求核对截止时间。系统需要辨认情绪线索与交流目的，而不是看见负面词就统一安慰。适当回应应承认已表达的处境，不添造经历或诊断；用户需要信息时提供有依据的信息，需要被倾听时不急着替其决定。下面方法研究普通支持性交谈，不能把主观喜好提高解释为专业处置效果。

\subsubsection{情绪条件生成与MoEL}
\label{sec:nlp-dialogue-moel}

基线可以先从历史预测情绪标签$e$，再以$(C,e)$生成回复。训练需要情绪标注与参考回复，可分别用分类交叉熵和词元负对数似然。缺点是硬选一个标签可能丢失混合感受；把真实标签提供给测试模型又会高估实际情绪识别能力。

MoEL由情绪跟踪器、多个情绪回应专家和一个汇总解码器组成。Transformer编码器用词、位置与说话者角色嵌入读取历史，查询位置表示$\bm{q}$对情绪键$\bm{k}_e$做softmax，得到$p_e$。每个专家是在同一历史和回复前缀上计算的独立Transformer解码层，输出$\bm{V}_e$；共享专家输出$\bm{V}_0$，不依赖某一种情绪。组合为
\begin{equation}
 \bm{V}_{\mathrm{mix}}=\bm{V}_0+\sum_e p_e\bm{V}_e,
 \qquad
 L=\alpha[-\log p_{e^*}]+\beta[-\log p(y^*\mid C)].
 \label{eq:nlp-dialogue-moel}
\end{equation}
汇总解码器进一步变换混合表示并输出词元分布。这里混合的是隐表示，不是多条完整回复的字符串，也不等价于词元概率的直接加权平均\scite{nlp-dialogue-lin2019}

例如“我第一次参加读书会，很期待”与“我第一次参加读书会，怕听不懂”共享事件而情绪不同。压缩成期待、焦虑两个教学专家时，前一句可给权重$(0.8,0.2)$，后一句为$(0.2,0.8)$。若某步$\bm{V}_0=(0.5,0.5)$、$\bm{V}_{\text{期待}}=(1,0)$、$\bm{V}_{\text{焦虑}}=(0,1)$，两次混合分别为$(1.3,0.7)$、$(0.7,1.3)$，再由汇总层生成不同回应。原实验使用更多情绪标签，这个二维例子只检查组合机制。

前一句可以回应期待并问感兴趣的话题，后一句可以承认担心并询问最担心哪类内容，不能编造“大家都会很耐心”这类未获依据的保证。若用户接着说“先别安慰，我只想知道地点”，情绪分布可能仍偏焦虑，但当前目的要求回答地点。固定历史和预算分别比较情绪识别准确率、回复相关性与人工感受到的理解程度；增加情绪监督不等于自动掌握用户没有说出的经历。

\subsubsection{ESConv支持策略条件生成}
\label{sec:nlp-dialogue-esconv}

ESConv把支持行为标成比情绪更接近回复职责的策略，例如提问、复述、反映感受、肯定与安抚、提供建议或信息。数据中的交流者依据指南选择策略并组织回应，模型由此学习“接下来怎样回应”与“回应具体怎样措辞”的联系。策略标签不是证明其对任何用户都合适的规则。

原论文以预训练对话模型为骨干，将策略特殊词元放在参考回复之前，以展平历史为条件，训练
\begin{equation}
 p(s,y\mid C)=p(s\mid C)\prod_i p(y_i\mid C,s,y_{<i}).
 \label{eq:nlp-dialogue-esconv}
\end{equation}
联合版本先预测或采样策略，再生成回复；给定正确策略的版本测表达能力，随机策略版本检查策略选择的贡献，不带策略的版本直接生成回复。三者输入信息不同，报告中必须分开，不能把正确策略实验称为端到端结果（见文献说明）。

以同一事件构造三种目的。“错过报名让我很失落，只想说说”宜先反映感受，如“你原本很期待参加，错过后确实有些失落”；“报名什么时候截止”宜提供通知中的精确信息；“能帮我想想怎样避免下次再错过吗”才适合提出与目标相符的提醒建议。假设后一句策略分布在反映感受、信息、建议上为$(0.2,0.1,0.7)$，选建议后生成“可以在确认下一次截止时间后设一个提前提醒”。恢复时既检查策略标签，也检查文字是否真的执行该策略，而不是输出建议标签却只说“别难过”。

支持策略中出现“自我披露”不授权系统虚构自己的真实经历；教学角色若有指定经历，应显式遵守角色条件，普通助手不能凭训练样本声称亲历。策略适合度、无依据推断、用户主观感受和实际帮助需分别记录。原研究的普通情绪支持交流及其参与者评价，不能推广为临床疗效或专业建议有效性；明显偏离需求的温暖语言也应计错。

\section{协作式对话}
\label{sec:nlp-dialogue-collaboration}

用户说“把通知写得简单些”，并没有完整规定交付物。“简单”可能指篇幅、阅读难度或语气；系统直接改短有时合适，有时删掉了用户最关心的条件。协作式对话把要求澄清、候选比较、反馈落实和最终验收连接起来。第\ref{chap:nlp-text-generation}章负责怎样生成与编辑，此处研究交流如何决定要生成什么，以及产物是否真正落实了当前要求。

\subsection{目标澄清}
\label{sec:nlp-dialogue-clarification}

澄清的输入是含混请求、已有材料和已确认条件，输出是有理由优先询问的问题，而不是越问越完整的问卷。对“简单些”，可以直接问希望多短，也可给两个不同版本让用户比较；哪种方式有效，取决于误解的代价、问题是否好回答及用户愿意投入多少。已确定的活动事实不应重复询问，尚无依据的业务事实则不能靠澄清措辞偷换成已知。

\subsubsection{神经EVPI澄清问题排序}
\label{sec:nlp-dialogue-evpi}

完全信息的期望价值（expected value of perfect information, EVPI）启发的神经排序模型估计补入回答之后，原请求是否更便于解决。原论文从相似论坛帖子取得澄清问题与对应回答候选，用问答表示模型评估某个回答与当前问题是否相配，用效用模型评估补充该回答的价值，再对候选问题排序。它不是直接生成任意澄清回答的语言模型，也没有把熵下降定义为效用\scite{nlp-dialogue-rao2018}

具体而言，词汇检索取得相似帖子，原设置使用十组问题及由后续编辑或作者回复取得的回答。设当前帖子为$p$，问题候选为$q_i$，检索到的成对候选为$(q_j,a_j)$。LSTM编码帖子和问题，前馈网络$F_{\mathrm{ans}}$输出一个回答表示；候选回答的平均词向量为$\widehat{\bm{a}}_j$。原文用余弦距离$d_{i,j}=1-\cos(F_{\mathrm{ans}}(p,q_i),\widehat{\bm{a}}_j)$及问题相似度形成回答权重。将其作为排序分数而非校准概率，可写成
\begin{equation}
\begin{aligned}
 w_{i,j}&=\exp(-d_{i,j})\,
            \cos(\widehat{\bm{q}}_i,\widehat{\bm{q}}_j),\\
 S(q_i;p)&=\sum_j w_{i,j}U(p,q_j,a_j).
\end{aligned}
\label{eq:nlp-dialogue-evpi}
\end{equation}
这里$U$由帖子、候选问题及候选回答的LSTM表示送入前馈网络，再经sigmoid得到。原式的候选得分未显式归一化为总和1，本书也不擅自补一个softmax后冒称原算法；教学数例限制问题相似度为非负，避免把带符号的相似度称作概率。

回答表示损失使预测向量接近原配对回答，并按问题相似程度接近其他候选回答。效用训练把原始帖子—问题—回答三元组作为正例，检索到的其他组合作为负例，以二分类损失训练；这些负例有噪声，因为另一个问题也可能有用。两个目标联合更新编码器与网络，监督来自论坛交互记录而不是逐问题的真实世界效用测量。

对“把通知写得简单些”，取$q_1$为“需要限制篇幅吗”，$q_2$为“是否面向首次参加者”；候选回答$a_1$为“限制在80字左右”，$a_2$为“面向首次参加者”。设两个问题的交叉余弦相似度为0.5，对角相似度为1，距离矩阵为
$\left[\begin{smallmatrix}0.2231&0.9163\\0.5108&0.5108\end{smallmatrix}\right]$，
得到权重近似$(0.8,0.2)$和$(0.3,0.6)$。再给定效用$U_1=0.7,U_2=0.9$，则$S(q_1)=0.74$、$S(q_2)=0.75$，排序先问受众。第二行权重和为0.9，未归一化；这里按原评分形式排序。

恢复输出时展示被选问题，等待真实用户回答，再更新目标；不能把候选回答$a_2$直接写成用户已经确认的事实。候选中没有询问语气的问题时，排序器无法选出它；效用模型的分数也没有扣除用户负担。原任务是候选澄清问题排序，本书将其接口移到通知协作属于教学迁移，不能据此声称已验证完整协作收益。

\subsubsection{规则询问与信息增益基线}
\label{sec:nlp-dialogue-information-gain}

规则基线可以先查必需条件：用途未定则问用途，两个硬要求冲突则先澄清冲突，其余采用声明的默认值并交付初稿。另一个可手算基线把未知目标表示为有限假设，询问能减少多少不确定性，再扣去交互成本。它需要显式给出先验、可能回答和成本，与上一小节学习到的候选效用不是同一个模型。

设$G$为目标假设，$A_q$为对问题$q$的回答，使用以2为底的熵：
\begin{equation}
 J(q)=H(G)-\sum_a p(a\mid q)H(G\mid A_q=a)-c(q).
 \label{eq:nlp-dialogue-information-gain}
\end{equation}
回答概率由假设先验与回答模型求得，成本$c(q)$换算为与信息收益可比较的约定单位；这种换算是教学设计，不是普适的“一轮值多少比特”。若回答可能误解或拒答，必须在$p(a\mid q,G)$中加入这些结果。

为涉及篇幅、难度和语气，规定四个互斥的教学目标：只缩短篇幅$g_1$，只降低难度$g_2$，只放松语气$g_3$，同时缩短并降低难度$g_4$，先验为$(0.4,0.3,0.1,0.2)$；未提到的属性保持原样。这不是所有可能目标的穷举。问题$q_L$问“是否主要需要更短”，把$g_1,g_4$归入“是”，回答概率为$(0.6,0.4)$；问题$q_D$问“是否需要降低阅读难度”，把$g_2,g_4$归入“是”，概率为$(0.5,0.5)$。此例假设回答确定且诚实。

先验熵为1.846比特。$q_L$的“是”后验在$g_1,g_4$上为$(2/3,1/3)$，“否”后验在$g_2,g_3$上为$(3/4,1/4)$，预期剩余熵为$0.6H(2/3,1/3)+0.4H(3/4,1/4)=0.875$，信息增益0.971。$q_D$的两组后验分别为$(0.6,0.4)$与$(0.8,0.2)$，预期剩余熵0.846，信息增益1。若$c(q_L)=0.10,c(q_D)=0.25$，则净分数分别为0.871和0.750，优先问篇幅。

直接交付最高先验目标$g_1$，在这个简化模型中完全猜对的概率为0.4；问$q_L$后分别选各分支最大后验目标，总命中概率为$0.4+0.3=0.7$。多问一轮确有潜在收益，却还未确定具体字数和难度等级。熵减只评价当前假设空间，可能与产物可用性不一致；最终仍需按真实交付质量和用户投入比较直接交付、给选项和先询问，不能把较高$J(q)$当成用户体验证明。

\subsection{候选讨论}
\label{sec:nlp-dialogue-candidates}

用户有时难以抽象说明语气，却能看出哪个版本适合。候选讨论要让差异可辨，同时保证比较不建立在一份正确、一份明显漏事实的伪选择上。点赞、选中某个版本或说“挺好”，都应按当时问题范围解释，不能自动视为所有细节已验收。

\subsubsection{约束覆盖与差异候选选择}
\label{sec:nlp-dialogue-diverse-candidates}

固定任务为发布当前活动简讯，已确认事实包括活动名、当前时间地点、校内参加条件和校外严格截止条件，不要求复述旧日期。生成器用这些条件产生候选；训练式生成可用条件—文本对，提示式生成固定材料与示例，规则筛选则检查字段覆盖。只有通过硬条件的候选才按篇幅、语气与组织方式比较，避免用多样性补偿事实缺失。

例如版本甲为“机器学习读书会将于2025年4月18日（周五）14:00在海棠楼201举行。校内师生可直接参加，校外来宾须在4月16日18:00前提交报名信息。”版本乙为“机器学习读书会：2025年4月18日14:00，海棠楼201。校内师生可直接参加；校外来宾须在4月16日18:00前提交报名信息。”乙压缩了句法和星期重复信息，保留相同条件，差异可以说明为完整叙述与紧凑简讯。

候选较多时可用第\ref{chap:nlp-text-generation}章的MMR式重排，以相关性$R(y)$和与已选集合$S$的相似度计算
\begin{equation}
 s(y)=\lambda R(y)-(1-\lambda)\max_{z\in S}\operatorname{sim}(y,z).
 \label{eq:nlp-dialogue-mmr}
\end{equation}
假设已选甲，$\lambda=0.7$，乙的相关性0.9、与甲相似度0.6，另一个几乎同文的候选丙为0.95、0.95，则乙得分0.45、丙0.38，选择差异更明显的乙。任何候选若漏掉18:00或“前”，应在硬条件过滤时淘汰，不能凭较低相似度入选。

用户说“第二个更短，但保留地点”，共同状态更新为选择乙的紧凑形式，并显式确认地点不可删。乙本来已经含地点，因此无需为了显示“已修改”再改一次地点；篇幅是否还要进一步缩短、是否省略年份仍未确认。比较一个结果、随机多个结果和差异候选时，要固定生成预算与展示数量，并计入用户阅读负担。MMR只选择可比较文本，不会自动识别用户赞成了哪些属性。

\subsubsection{Bradley--Terry成对偏好建模}
\label{sec:nlp-dialogue-bradley-terry}

成对选择积累后，可以学习相对偏好。Bradley--Terry模型给每个候选一个分数$s_i$，把得分差映射为选择概率：
\begin{equation}
\begin{aligned}
 p(i\succ j)&=\frac{\exp s_i}{\exp s_i+\exp s_j}
             =\sigma(s_i-s_j),\\
 L&=-\sum_{(i,j)}\bigl[n_{i,j}\log p(i\succ j)
               +n_{j,i}\log p(j\succ i)\bigr].
\end{aligned}
 \label{eq:nlp-dialogue-bt}
\end{equation}
$n_{i,j}$为在固定问题、展示协议和用户条件下选择$i$胜过$j$的次数。分数整体加同一常数不改变概率，估计时需固定一个分数为零或令分数和为零；比较图若不连通，不同连通分量间的次序没有数据依据。参数由比较似然估计，推广到新候选时可用条件特征网络输出分数，已确认硬要求仍保留为条件。

以同一用户在相同目标下对甲乙的四次教学比较为例，假设乙胜3次、甲胜1次，且暂按独立同分布选择建模。最大似然估计为$p(\text{乙}\succ\text{甲})=3/4$；固定$s_{\text{甲}}=0$，得到$s_{\text{乙}}=\log3=1.099$，负对数似然为$-3\log(3/4)-\log(1/4)=2.249$。只出现单边胜出的数据会使无正则估计趋向无穷，实际需加入正则或先验，而不是报告无限偏好。

本例仅两候选，所以排序为乙在前。若有平局、不确定或“选乙但恢复日期”，应保留独立标签及修改文本；可以按预先约定排除平局或采用支持平局的扩展，不能悄悄当成一次乙获胜。偏好训练的一般思想可联系人类比较学习\scite{christiano2017}

重复询问同一用户也可能受顺序、记忆和疲劳影响，独立性是假设而非事实。分数描述某种任务条件下的相对选择，不证明乙每一条事实都对，更不表示用户接受了省略地点。实际交付还需检查已确认约束和最新反馈；聚合其他用户偏好也不能覆盖当前用户的明确要求。

\subsection{互动修订}
\label{sec:nlp-dialogue-revision}

候选选定后，交流的结果应体现在产物里。系统保存版本号、当前已确认条件和未决问题，将每次反馈绑定到其所针对的版本，避免用户评论旧稿时误改新稿。停止依据是用户验收或预先约定的任务检查，而不是已经回复了一句“已修改”。

\subsubsection{用户反馈条件编辑}
\label{sec:nlp-dialogue-feedback-edit}

用户反馈条件编辑先把意见转成目标状态的增量，再调用局部编辑或条件改写，最后检查实际差异。输入包括原稿、事实材料、反馈及其作用范围；训练式模型可用原稿—反馈—修订稿三元组，提示式编辑固定同一接口，规则也可直接替换已定位字段。第\ref{chap:nlp-text-generation}章讨论文本编辑算法，此处增加反馈来源、确认状态和版本一致性。

设有一份明确的教学旧稿，两处写了“4月17日”，地点仍为海棠楼201。用户说“改成4月18日，地点不变”，其依据是共同通知的已确认改期，而不是用户凭空修改真实活动。状态变化只允许覆盖当前活动日期，编辑器将标题或正文中两处当前日期同步替换；若保留“原定4月17日”的历史说明，那一处不应改，因为它表示旧计划属性。

用压缩词元序列检查最小例子：旧稿为“读书会／4月17日／海棠楼201／请于／4月17日／到场”，目标为“读书会／4月18日／海棠楼201／请于／4月18日／到场”。单位代价的Levenshtein对齐需要两次替换，其他词元保持。若结果只替换一次，必要修改完成率为$1/2$；若又把地点改成另一楼，虽然两日期均修复，却多出一次不被授权的变化。差异显示与字段核验共同恢复这些错误。

有歧义的“改到周五”需先绑定日期属性和年份，不能改报名截止；用户反馈若与权威事实冲突，要说明冲突并澄清是在编辑假设通知还是修订真实材料。多参与者协作还需规定谁可确认事实、谁只负责措辞，并保留来源。多个发言身份只是工作分工，不因此成为多智能体强化学习；最终需检查共同产物、必要变化、额外变化、轮次和用户投入。

\subsubsection{Self-Refine辅助自反馈修订}
\label{sec:nlp-dialogue-self-refine}

Self-Refine让同一个已训练语言模型依次承担初稿生成、反馈和修订。原方法不额外训练专门的反馈器，也不把这一循环称为强化学习；能力来自基础模型及三类任务提示中的说明或少量示例。设请求与事实材料为$x$，初稿为$y_0$，反馈为$f_t$，则过程为
\begin{equation}
\begin{aligned}
 y_0&=M(p_{\mathrm{gen}},x),\\
 f_t&=M(p_{\mathrm{feedback}},x,y_t),\\
 y_{t+1}&=M(p_{\mathrm{refine}},x,y_0,f_0,\ldots,y_t,f_t).
\end{aligned}
\label{eq:nlp-dialogue-self-refine}
\end{equation}
原文在迭代修订提示中保留此前输出与反馈，直到固定轮数或反馈中的停止指示满足条件。停止指示仍是模型输出，不是独立的事实核验；具体机制见文献说明。

在通知任务中，假设初稿把要求写成“校外来宾在4月16日前报名”。反馈提示要求“对照给定材料，定位被删去的时间与限定词，只指出可执行修改”，可得到“补回18:00及严格的前，并保持活动日期地点不变”。修订提示连同原材料和反馈生成“校外来宾须在4月16日18:00前提交报名信息”，其他字段保持。按本书规则检查器核对该句的截止时刻、严格不等号、主体与动作，才确认本次修改没有改变条件。

本例使用一次初稿调用、一次反馈调用和一次修订调用；若已有用户初稿，生成初稿的成本可另记，不能仍把整个流程算成一次回答。可以规定最多两轮自修订，超过后交还用户或说明未解决项。与直接按人类反馈编辑相比，Self-Refine减少了某些检查对用户输入的依赖，却增加模型计算，也可能重复自己的错误。

若模型自反馈提出“删除地点以便更短”，而用户已确认地点必留，编辑器必须拒绝这条建议；自反馈没有用户指令的优先级。反馈说“现在完全正确”也不能代替引用支持或用户验收。协作流程中可把它作为内部检查手段，分别记录用户轮数与自循环调用数，用最终产物的独立检查衡量收益，而非用模型自评文字衡量。

\section{问答与对话的综合比较}
\label{sec:nlp-dialogue-comparison}

问答与对话可以共享模型，却应按不同完成条件比较。固定问答资料版本、可见片段与总词元、检索和生成预算；固定对话服务快照、接口权限、真实用户或模拟器协议、最大轮数与中断条件。允许多种正确表达、等价证据组合和合法交流路径时，应在评分前写出等价规则，而不是见到某模型输出后才扩展。

问答的规范化准确匹配检查答案是否与某一允许参考一致；词元F1则由预测与参考的多重集合交集计算精确率和召回率。例如参考“4月18日／14:00”，预测只有“4月18日”，精确率1、召回率$1/2$、F1为$2/3$，准确匹配为0。若时间与报名条件都需回答，而召回只有片段A，按两项所需证据计的召回率为$1/2$。若生成了时间和截止两个正确陈述，却两句都引用A，引用支持率仍只有$1/2$。答案、覆盖和引用三个分数不能互相替代。

状态评价把每轮有效的槽位—值对作为集合，报告槽位F1和联合目标准确率。表\ref{tab:nlp-dialogue-metrics}沿用三轮主例，假设某系统前两轮正确，第三轮错误保留4/17。为避免空集合约定隐藏结果，本例将两边都空的轮次单独记为正确；非空槽位采用跨轮汇总的微平均。

\begin{table}[htbp]
\centering
\begin{tabular}{cllc}
\toprule
轮次 & 正确状态 & 假设预测状态 & 全状态正确\\
\midrule
1 & 日期、人数均缺失 & 日期、人数均缺失 & 是\\
2 & 4/17，2人 & 4/17，2人 & 是\\
3 & 4/18，2人 & 4/17，2人 & 否\\
\bottomrule
\end{tabular}
\caption{局部状态分数与完整任务的差别。活动编号和身份为固定上下文，不计入可修改槽位；第三轮旧日期使查询继续为空。}
\label{tab:nlp-dialogue-metrics}
\end{table}

三轮共4个正确正值、4个预测正值，匹配3个，所以微平均精确率、召回率和F1均为$3/4$，联合目标准确率为$2/3$。第三轮唯一需要变化的日期没有改对，其变化项召回率为0；人数保持正确不抵消更正失败。这个系统可以每轮都生成礼貌回复，却没有查到用户要核对的报名，完整任务成功为0。正确主线则在第三轮取得匹配记录并检查提交时点，以当前目标定义成功，不以“调用次数达到两次”定义成功。

完整交互还应报告用户补充或更正次数、总轮数、等待时长、检索与模型调用成本。故障造成的重试和澄清目标造成的新轮次应区分；自反馈调用也不能记为用户轮次。若环境允许用户同时修改共享状态，必须在轨迹中保存这些观察，不能继续用静态快照评分。

比较两个系统时，以相同问题或同一初始对话任务配对。假设四个教学任务上的最终成功分别为A的$(1,1,0,0)$和B的$(1,0,1,0)$，两者平均成功率均为$1/2$，逐任务差为$(0,1,-1,0)$。它们修复的是不同任务，不能据平均值相同就说失败机制相同。样本足够时可按完整任务配对重采样计算差值区间，多轮对话不能把各轮当成独立任务来虚增样本量。

人工选择、Bradley--Terry聚合、指定多轮评分标准的模型裁判与实际任务证据分别记录。裁判需要看到被评价方法实际可用的资料与轨迹，检查顺序偏差，并用独立人工或可执行任务依据校验；它的总分不是接口真值。诊断时依次定位资料是否取得、本轮是否理解、状态是否更新、动作是否合规、回复是否有据、环境目标是否成立。给正确证据或正确状态的诊断实验有助于隔离环节，但不能替代端到端比较。

\section*{文献说明}
\label{sec:nlp-dialogue-literature}

以下集中列出较长或相邻出现的完整文献，便于在不挤占边注的情况下核对方法来源。数值例、通知资料、只读报名接口及规则恢复协议均为本章教学构造，不是这些论文的实验结果。

\begin{fullwidth}
\raggedright
\fullcite{yang2018hotpotqa}\par
\fullcite{nlp-karpukhin2020dpr}\par
\fullcite{lewis2020rag}\par
\fullcite{nlp-dialogue-izacard2021}\par
\fullcite{nlp-dialogue-herzig2020}\par
\fullcite{nlp-dialogue-zhu2021}\par
\fullcite{nlp-dialogue-chen2021}\par
\end{fullwidth}

上述资料支持多跳证据任务、DPR对比检索、RAG边际化、TAPAS表格选择与聚合、TAGOP标签和算子预测，以及FinQANet限定程序生成。引用标记核验和跨工具权限控制是本书额外规定，不是模型概率本身的保证。

\begin{fullwidth}
\raggedright
\fullcite{min-etal-2023-factscore}\par
\fullcite{nlp-parikh2016decomposable}\par
\fullcite{budzianowski2018multiwoz}\par
\fullcite{rastogi2020sgd}\par
\fullcite{nlp-dialogue-chen2019}\par
\fullcite{nlp-dialogue-hosseini2020}\par
\fullcite{nlp-intro-barres2025}\par
\end{fullwidth}

这里分别对应原子事实精确性、可分解注意力、JointBERT轮次标签、SimpleTOD序列化任务接口及双控制环境评价。支持性判断、状态正确与环境成功仍按本章相应协议分别验收。

\begin{fullwidth}
\raggedright
\fullcite{nlp-intro-dinan2019}\par
\fullcite{nlp-intro-zhang2018persona}\par
\fullcite{nlp-dialogue-liu2021}\par
\fullcite{nlp-dialogue-madaan2023}\par
\end{fullwidth}

知识交流和人物条件的模型实现以原文的方法节为准；Self-Refine使用原文的生成、反馈、修订提示接口。用户记忆规则、通知差异候选和信息增益基线是为解释交流过程给出的可复算实现，不混称为这些论文的原始算法。

\section{本章小结}
\label{sec:nlp-dialogue-summary}

从回答一句话扩展到可靠交互，需要保存的不只是更长的上下文。证据问答将问题连接到资料、跨度、计算或生成结果，核验再把答案拆成可以分别支持、反驳或限定的陈述；任务型对话维护当前有效的需求，把本轮理解、状态变化、行为和真实结果分开，才能在查询失败或用户更正时恢复。开放交流还要区分知识依据、人物条件和支持需求，协作交流则把澄清、选择和反馈落实到实际产物。

这些任务都能使用检索、注意力与条件生成，但正确的完成条件来自任务本身。一个高概率答案可能缺证据，一条合法状态可能已经过期，一次格式正确的调用可能尚未执行，一个被偏好的回复也可能没有解决问题。信息不足时保留未知、提出必要澄清或只交付已获支持的部分，比用流畅表达掩盖缺口更可检查。下一章进一步把已明确的语言要求转为解答、程序和环境行动，并以相应验证对象判断任务是否完成。
